\documentclass[professional,privacy,print,win,long,STEM,AutoFakeBold]{def/buaa}
\refcolor{on} 
\usepackage{geometry}
\usepackage{graphicx}
\usepackage{booktabs}
\usepackage{multirow}
\usepackage{amsmath}
\usepackage{amssymb}
\usepackage{url}
\usepackage{tocloft}
\usepackage{titlesec}
\usepackage{hyperref}
\usepackage{float}
\usepackage{xeCJK}
\usepackage{fontspec}
\geometry{left=3cm, right=2.5cm, top=2.5cm, bottom=2.5cm}
\graphicspath{{./pic/}} % 设置图片搜索路径
% 目录格式优化
\setlength{\cftbeforesecskip}{4pt}
\renewcommand{\contentsname}{\centerline{目\quad 录}}



%%=================================================================
% 论文题目及副标题-{中文}{英文} 注意：论文题目应严格控制在25个汉字（符）以内 
\Title{基于LLM微调的银行金融财务推理领域的应用优化}
      {Optimization of Banking Financial Reasoning Based on Fine-Tuning of LLMs}
% \Subtitle{版本 \BUAAThesisVer{}}{Version \BUAAThesisVer{}}

\Branch{工学}

\Degree{Master of Software Engineering}

% 院系-{中文}{英文}，专业，研究方向，一级学科，学科方向（学术型学位）
\Department{软件学院}{School of Software}
\Major{软件工程}
\Feild{软件工程}
\Discipline{软件工程}
\Direction{软件工程}

% 导师信息-{中文名}{英文名}{职称}
\Tutor{周号益}{Haoyi Zhou}{副教授}

% 学生姓名-{中文名}{英文名}
\Author{叶晖}{Yehui}
% 学生学号
\StudentID{ZF2321310}

% 中图分类号
\CLC{待补充}

% 时间节点-{月}{日}{年}
\DateEnroll{12}{01}{2025}
\DateGraduate{07}{01}{2026}
\DateSubmit{12}{01}{2025}
\DateDefence{12}{01}{2025}
%%=================================================================
% 摘要-{中文}{英文}
\Abstract{%
  随着金融行业数字化转型的深入，银行业务对智能化财务分析与决策的需求日益迫切。传统金融文本处理技术在处理复杂财务推理任务时，普遍面临逻辑连贯性不足、专业术语误判以及跨模态数据融合能力弱等瓶颈。现有大语言模型虽在通用场景展现潜力，但直接应用于金融领域时仍存在数值推理偏差、可解释性不足等系统性缺陷。本研究聚焦金融财务场景，探索基于大语言模型微调的财务推理优化方法，旨在构建兼顾专业性、准确性与合规性的智能分析框架。

  研究以金融财务领域的数据处理与可解释推理为核心，通过领域知识注入与模型架构优化，系统解决文本语义理解偏差、推理逻辑断裂等关键问题。设计三阶段递进训练架构，首先训练基于Token分类的Process Reward Model（PRM）实现步骤级过程监督（步骤准确率达95\%），进而基于PRM构建大规模偏好数据集（平均区分度0.69），通过DPO偏好对齐训练优化金融推理策略模型，在留出集上实现99\%以上的偏好准确率；端到端三维度评估显示，策略模型的最终结论正确率较基座提升11.75个百分点（相对25.5\%），验证了偏好对齐对金融推理正确性的实际提升。研究预期形成面向银行业务的智能财务分析系统，显著提升财务文本理解与数值推理的准确性，为信贷审批、风险预警等核心场景提供高效、可靠的技术支撑。

  }{
  With the deepening digital transformation of the financial industry, the demand for intelligent financial analysis and decision-making in banking is increasingly urgent. Traditional financial text processing techniques often face bottlenecks when handling complex financial reasoning tasks, including insufficient logical coherence, misinterpretation of domain-specific terminology, and weak cross-modal data integration. Although existing large language models (LLMs) show promise in general scenarios, their direct application in the financial domain still exhibits systematic limitations, such as numerical reasoning bias and lack of interpretability. This study focuses on financial and accounting scenarios, exploring LLM-based fine-tuning methods for optimized financial reasoning, aiming to construct an intelligent analysis framework that balances professionalism, accuracy, and compliance.

  The research centers on data processing and interpretable reasoning in the financial domain, systematically addressing critical issues such as semantic understanding deviation and reasoning logic disruption through domain knowledge injection and model architecture optimization. We design a three-stage progressive training architecture: first, a Token-classification-based Process Reward Model (PRM) is trained to provide step-level process supervision (step accuracy 95\%); then, a large-scale preference dataset (average margin 0.69) is constructed based on PRM scoring, and DPO-based preference alignment training is applied to optimize the financial reasoning policy model, achieving over 99\% preference accuracy on held-out sets. End-to-end three-dimensional evaluation shows that the final conclusion accuracy of the policy model improves by 11.75 percentage points (a relative 25.5\%) over the base model, verifying the practical gain of preference alignment on financial reasoning correctness. This study is expected to result in an intelligent financial analysis system for banking applications, significantly improving the accuracy of financial text understanding and numerical reasoning, and providing efficient and reliable technical support for core scenarios such as credit approval and risk early warning.
}
% 关键字-{中文}{英文}
\Keyword{大语言模型，金融财务推理，过程奖励模型，参数高效微调，数据合成，DPO}{Large Language Models, Financial Reasoning, Process Reward Model, Parameter-Efficient Fine-Tuning, Data Synthesis, DPO}

% 图表目录
\Listfigtab{on} % 启用: on[默认]; 关闭: off;

%----------------- 文档开始 -----------------%
\begin{document}
%----------------- 正文 -----------------%
\newpage
\chapter{绪论}
\section{课题背景}
本课题的提出，主要源于金融行业在数字化转型过程中面临的现实压力。这种压力来自三个方面：国家政策的引导、人工智能技术的快速进步，以及金融行业自身长期存在的实际问题。

从政策角度看，银保监会在《金融科技发展规划（2022–2025年）》中明确提出，金融机构要“持续优化风险控制模型和策略”，并加快构建基于人工智能的自动化决策系统。这一要求意味着，过去那种主要依靠人工经验来分析金融文本、进行数值推理的做法，已经无法满足当前监管对效率和可追溯性的硬性标准。

从技术角度看，以Transformer架构为基础的大语言模型近年来取得了重要突破。这些模型在理解长文本、进行多步逻辑推理等方面的能力，明显强于传统方法。例如，程序化推理链（Programmatic CoT）和过程奖励模型（Process Reward Model, PRM）等新技术，为解决金融领域复杂的推理任务提供了可行方案。这为金融业务的智能化升级打开了新的窗口。

从行业实践角度看，财务推理任务如财务报表分析、投资策略制定和合规评估，要求一定程度的精确性、事实性和逻辑连贯性，突破现有模型的极限。目前金融机构在财务分析、风险评估等关键业务中，依然存在明显的技术短板。特别是在处理非结构化文本（比如财报说明）和结构化数据（比如财务报表）的联合推理任务时，现有系统常常出现逻辑不连贯、专业术语理解错误等问题。在CFLUE、FinQA等金融专业测试集上，传统深度学习模型的表现长期停滞，暴露出数值推理链条断裂、多源信息无法有效融合等系统性缺陷。这些问题与金融行业日益追求精细化、实时化决策的趋势形成了突出矛盾。

因此，如何借助大语言模型的新能力，构建更可靠、更专业的金融推理系统，成为一项紧迫且具有现实意义的研究课题。

\section{课题意义}
本课题立足金融科技发展前沿，针对银行场景下的财务智能分析需求，探索大语言模型在金融专业领域的深度应用优化路径。

在理论层面，现有研究普遍存在跨模态知识割裂问题，传统自然语言处理技术难以有效融合财务报表的数值特征与文本语义信息。本研究拟构建金融财务领域适配的语义-数值联合推理框架，突破传统方法在金融文本理解与数值计算间的协同障碍，为金融智能系统提供可解释的推理范式。通过建立面向银行场景的知识表征与推理机制，有望增强金融认知计算领域财务推理的技术能力，完善智能财务分析的理论体系架构。

在实践应用层面，商业银行日常运营面临海量财务文档处理压力。基于大语言模型的智能分析系统可显著提升金融数据处理效率与决策质量：在信贷审批场景，通过自动化财务指标计算与风险特征提取，可将传统人工处理复核效率提升40\%以上；在反欺诈检测领域，融合多源异构数据的动态推理机制能够有效识别隐蔽的财务舞弊模式；在合规审计环节，程序化推理链技术可建立完整的过程追溯路径，满足巴塞尔协议III的监管要求。本研究成果将为银行业数字化转型提供关键技术支撑，对推动智能风控体系建设、降低金融机构运营成本具有重要现实价值。


\section{国内外研究现状}

\subsection{国外研究现状}

\subsubsection{金融财务推理领域进展}

近年来，金融领域专用大语言模型在精度与效率层面实现显著突破。以Fino1\upcite{qianFino1TransferabilityReasoning2025}、DANA\upcite{luongDANADomainAwareNeurosymbolic2024}为代表的领域优化模型，通过融合神经符号计算与动态知识注入技术，在财务报表分析、风险预警等任务中展现出超越通用大模型的性能优势。例如Fino1模型通过程序约束式推理机制，在XBRL-Math财务公式解析任务中达到82.22\%的正确率，验证领域专用架构的技术优越性。与此同时，轻量化训练策略（如LoRA\upcite{huLoRALowRankAdaptation2021}、MoE\upcite{jiangMixtralExperts2024}）的普及使模型部署成本大幅降低，为中小金融机构的技术落地提供可行性路径。

国际学术界在金融大模型领域取得重要突破，围绕基础模型构建、多模态推理和程序化生成三个维度形成技术体系。在基础模型层面，Wu等（2023）研发的BloombergGPT\upcite{wuBloombergGPTLargeLanguage2023}通过领域自适应预训练策略，在金融命名实体识别任务中展现出专业术语处理优势，但其结构化数据处理能力仍显不足，特别是在财务报表数值计算场景存在系统性偏差。

当前国际研究在动态知识更新和领域适配效率方面仍存在显著瓶颈。主流模型普遍采用全参数微调策略，导致训练成本居高不下，且难以适应金融市场的时效性要求。在可解释性层面，现有方法大多依赖后验分析技术，未能实现推理过程的内在约束，这在涉及监管合规的财务审计场景中形成实质性障碍。此外，针对长文档处理的注意力机制优化仍停留在通用文本层面，缺乏对财务报告特有结构（如附注与主表的勾稽关系）的针对性设计，导致关键信息提取效率低下。这些技术缺陷制约了大模型在银行核心业务场景的深度应用，亟待通过领域知识注入和推理架构创新实现突破。

\paragraph{参数规模与训练策略优化}

参数规模扩展已从单一维度转向计算效率优化\upcite{qinO1ReplicationJourney2024,grattafioriLlama3Herd2024}。Hoffmann等人\upcite{hoffmannTrainingComputeOptimalLarge2022}通过400组模型实验发现，计算最优训练要求模型参数规模与训练数据量等比例扩展（如70B参数的Chinchilla模型在4倍数据下，MMLU准确率达67.6\%）。Phi-4\upcite{abdinPhi4TechnicalReport2024}以14B参数证实小模型性能突破：通过反事实数据增强与训练课程优化，在HumanEval代码生成基准达到82.6\%准确率（超过部分百亿级模型）。Llama 3\upcite{grattafioriLlama3Herd2024}采用过度训练策略：405B参数模型训练运算量相比Llama 2最大模型的训练计算量增加近50倍，由此实现推理效率与下游任务性能平衡，其数学推理能力在GSM8K基准达到89.0\%。参数扩展边际效益递减现象驱动替代策略发展，如Snell等\upcite{snellScalingLLMTestTime2024}验证推断时间扩展的高效性——当推理计算量较低时（R=0.16），采用测试时计算优化的较小模型（PaLM 2-S*）可达到14倍参数规模模型的性能水平。这表明在特定场景下，测试时增加计算量可替代增加参数量。

\paragraph{多模态语义理解架构}

多模态处理正从简单图文对齐转向复杂结构化数据分析。Llama 3（Grattafiori等，2024）通过模块化适配器架构实现多模态扩展\upcite{grattafioriLlama3Herd2024}：在MMMU多模态推理验证集（900图像30学科任务）达到64.5\%准确率，较GPT-4V提升8.1pct，但略低于Claude 3.5 Sonnet的68.3\%。其图像编码器支持像素级图表解析，在ChartQA数值推理任务实现85.8\%结构识别准确率。非对称架构设计（独立视觉编码器+冻结语言主干）避免模态干扰，保持语言理解能力稳定。

\paragraph{领域轻量化适配技术}

领域适应技术呈现低开销与高精度融合趋势，Tulu3详细探讨了各种post-training方法\upcite{lambertTulu3Pushing2025}，包括监督微调（SFT）、直接偏好优化（DPO）和我们称之为可验证奖励强化学习（RLVR）的一种新方法。本研究评估了以下几个关键方法：

LoRA微调：在RoBERTa大型模型（355M参数）应用中，LoRA（Hu等，2021）仅需0.8M可调参数即达到全参数微调同等水平（MRPC任务F1 90.9±1.2）。GPT-3 175B适配时，LoRA将训练参数压缩4.7M/37.7M，预训练权重冻结策略保持基础能力完整（MNLI任务准确率波动＜0.5\%）\upcite{huLoRALowRankAdaptation2021,bidermanLoRALearnsLess2024}。

MoE架构：Mixtral通过专家路由机制优化长文本处理，在passkey检索任务中实现100\%准确率（32k上下文长度），且长文本困惑度随上下文增长单调下降\upcite{jiangMixtralExperts2024}。

检索增强：Borgeaud等\upcite{hoffmannTrainingComputeOptimalLarge2022}的检索式预训练将有效数据规模扩展10倍，使Codex代码生成HumanEval通过率提升27pct（34\%→61\%），验证外部知识整合对能力泛化的关键作用。

\paragraph{奖励模型和蒙特卡罗方法}

在大语言模型（LLM）的数学推理任务中，奖励模型被广泛用于提升推理的可靠性。目前主流的方法包括基于最终结果的奖励模型（Outcome Reward Model，ORM）和基于推理过程的奖励模型（Process Reward Model，PRM）。这两类模型的核心差异在于监督信号的粒度。

ORM仅对模型输出的最终答案进行判别。给定一个问题 $q$ 与其生成的解答 $x$，ORM 学习映射：
\[
p = \mathrm{ORM}(q, x) \in [0,1],
\]
其中 $p$ 表示最终答案正确的概率。训练目标通常为二分类交叉熵损失：
\[
\mathcal{L}_{\mathrm{ORM}} = -\frac{1}{N}\sum_{i=1}^N \left[ y_i \log p_i + (1 - y_i)\log(1 - p_i) \right],
\]
其中 $y_i \in \{0,1\}$ 为答案是否正确的标签。ORM 的优势在于标注成本低，可自动通过答案比对获取标签。然而，它无法识别中间步骤的逻辑错误。例如，模型可能通过错误推理偶然得到正确答案，ORM 会将其误判为高质量输出。这种“欺骗性路径”会误导强化学习或重排序过程，削弱模型的泛化能力 \upcite{lightmanLetsVerifyStep2023a}。

相比之下，PRM 对推理链中的每一步进行监督。设解答 $x$ 被划分为 $T$ 个步骤 $x = (x_1, x_2, \dots, x_T)$，PRM 对每个前缀 $x_{1:t}$ 和当前步 $x_t$ 进行评分：
\[
p_t = \mathrm{PRM}([q, x_{1:t-1}], x_t).
\]
PRM 的训练同样采用分类损失，但标签粒度细化到每一步。Lightman 等人 \upcite{lightmanLetsVerifyStep2023a} 通过人工标注构建了包含 80 万步标签的 PRM800K 数据集，并证明 PRM 在 MATH 基准上显著优于 ORM，其最佳模型解决了 78.2\% 的问题，而 ORM 仅为 72.4\%。

然而，PRM 面临两大挑战：一是高质量过程标注成本极高；二是“作弊”问题——即使每步得分高，最终答案仍可能错误。为解决标注瓶颈，近期研究转向自动化过程监督。Luo 等人 \upcite{luoImproveMathematicalReasoning2024} 提出 OmegaPRM，一种基于蒙特卡罗树搜索（MCTS）的自动化标注框架。其核心思想是通过 rollout 快速定位推理链中的首个错误。

具体而言，给定一个部分解答 $x_{1:t}$，模型从该状态生成 $k$ 条 rollout 路径，并统计其中能得出正确答案的比例：
\[
c_t = \frac{\text{正确 rollout 的数量}}{k}.
\]
该比例 $c_t$ 作为步骤 $t$ 的软标签。为提升效率，OmegaPRM 采用二分搜索策略定位首个 $c_t = 0$ 的位置，将时间复杂度从 $O(kT)$ 降至 $O(k \log T)$。在此基础上，该方法构建一棵状态-动作树，每个节点存储访问次数、蒙特卡罗估计值等统计量，并通过一种改进的 PUCT 公式选择最有价值的 rollout 进行探索。最终，OmegaPRM 在无需人工干预的情况下生成了 150 万条过程监督数据，显著提升了 Gemini Pro 和 Gemma2 模型的数学推理性能。

Cao 等人 \upcite{caoMoreBangBuck2025} 进一步提出基于熵驱动不确定性的 PRM（EDU-PRM），完全摒弃了规则化的步骤划分。他们观察到，在解码过程中，模型对下一个 token 的预测熵能有效指示逻辑转折点。高熵位置意味着模型处于不确定状态，往往是推理分支的关键节点。形式上，第 $v$ 个 token 的熵定义为：
\[
H_v = -\sum_{v} P_v \log(P_v + \epsilon),
\]
其中 $P_v$ 是语言模型输出的 token 概率分布，$\epsilon$ 为数值稳定项。EDU-PRM 将熵超过自适应阈值 $\tau(H)$ 的位置视为“不确定性锚点”，并以此动态分割推理步骤。

在锚点处，EDU 采样策略仅对 top-2 logits 进行分支，后续 token 采用贪婪解码，以平衡多样性与计算开销。每个生成的解答片段通过蒙特卡罗估计（MCE) 获得正确性标签。实验表明，仅使用 1.5\% 的训练数据，EDU-PRM 就能在 ProcessBench 上匹敌甚至超越 SOTA 模型。在推理阶段，EDU 采样策略不仅将准确率从 64.7\% 提升至 67.3\%，还将 token 消耗降低了 32\%，展现了极高的数据与计算效率。

OpenAI 的 o1 模型 \upcite{openaiOpenAIO1System2024} 虽未明确使用 PRM 术语，但其“ deliberative alignment”（审慎对齐）机制与 PRM 思想高度一致。o1 通过大规模强化学习，学会在输出前进行长链式思考，并在此过程中主动应用安全策略。这种内置的、显式的推理过程使其能更可靠地遵循内容规范，在 jailbreak 评测中显著优于 GPT-4o。这从侧面印证了过程监督对于提升模型鲁棒性和可控性的关键作用。

在开源领域，DeepSeek-R1\upcite{deepseek-aiDeepSeekR1IncentivizingReasoning2025}通过大规模强化学习训练，自然涌现出长链式推理能力，在多项推理基准上达到与o1可比的性能。其采用的GRPO方法\upcite{shaoDeepSeekMathPushingLimits2024}通过组内相对奖励替代价值函数，有效降低了RL训练的内存开销。DeepSeek-R1的成功表明，过程监督与RL的结合是实现强推理能力的可行路径，为本研究的PRM+DPO技术路线提供了重要佐证。

综上，从依赖人工标注的 PRM800K，到基于 MCTS 的自动化 OmegaPRM，再到利用内在不确定性的 EDU-PRM，过程奖励模型正朝着更高效、更可扩展的方向发展。蒙特卡罗方法作为连接最终结果与中间步骤的桥梁，是实现自动化过程监督的核心技术。未来的研究可进一步探索如何将过程监督与推理树搜索、动态阈值调整等技术结合，以构建更强大的通用推理智能体。


\subsubsection{国内研究现状}

现有研究在应对实际业务需求时仍面临多重挑战\upcite{LiMengXinAnLiTuiLiQuDongShuJuTeZhengDeQiYeCaiWuWeiJiYuCeYanJiu2024}。首先，高质量金融数据资源的稀缺性制约了模型性能提升，特别是在涉及企业敏感信息的场景中，数据获取与隐私保护的矛盾尚未有效解决。其次，模型推理效率与业务实时性要求之间存在显著差距，当前主流大模型部署要求高于中小银行承受能力，处理复杂财务问题的响应效率无法满足银行业务系统需求。这些瓶颈问题的突破将成为推动金融大模型实用化的关键所在。Xie等构建的FinBen\upcite{xieFinBenHolisticFinancial2024}是首个大规模开源金融LLM评估基准，涵盖36个数据集、24项任务、七个评估维度，其系统评测揭示了现有LLM在金融高级推理与复杂任务生成方面的普遍不足，进一步印证了上述挑战的系统性。

\paragraph{财务报告分析与欺诈检测}

该领域研究聚焦于复杂财务文本的结构化解析与风险特征识别。以Yang等人提出的FinGPT\upcite{yangFinGPTOpenSourceFinancial2023}框架为例，通过轻量化低秩适应技术（LoRA）构建动态知识图谱，实现对上市公司财报数据的语义解析特征增强。该框架在金融欺诈检测任务中表现突出，其采用三层异构数据处理机制：底层依托SWRL规则实现财务指标间的逻辑推理，中间层引入改进型金融情感极词表（如FinBERT的领域词典扩展方法），顶层结合本体知识库的动态更新算法，有效解决了虚假财务信息识别中对上下文语义敏感的问题。Zhang团队开发的XuanYuan 2.0模型\upcite{zhangXuanYuan20Large2023}在中文金融领域对话任务中展现出上下文适配的响应能力。

\paragraph{风险预测与信用评估}

研究主要突破传统结构化指标的局限，探索非结构化数据隐含风险的量化评估。在商业银行信用评估实践中，Liu等人提出的CBR混合推理引擎\upcite{LiMengXinAnLiTuiLiQuDongShuJuTeZhengDeQiYeCaiWuWeiJiYuCeYanJiu2024}研究突破传统结构化指标的局限，探索非结构化数据的风险量化评估：后端通过案例推理（CBR）系统对违约案例库进行相似性检索与结果修正。该方法使小微企业信用评估的AUC指标达到0.87。Li等开发的ICE-PIXIU系统\upcite{huNoLanguageIsland2024}通过构建中英双语指令数据集ICE-FIND和评估基准ICE-FLARE，基于InternLM模型微调出ICE-INTENT，在跨语言金融任务中实现了对GPT-4中文金融任务性能的超越，验证了翻译数据对模型泛化能力的提升作用。

\section{研究目标和研究内容}

\subsection{研究目标}

本课题面向银行金融场景下的智能化财务决策需求，针对现有大语言模型在财务文本理解、数值推理和合规生成等方面存在的关键瓶颈，重点突破跨模态语义融合与可解释推理技术难题。研究目标体系构建遵循“问题定位-技术突破-效能验证”的递进逻辑，具体从以下维度展开：

核心研究目标旨在建立具有金融领域适应性的智能推理框架。基于Wind/Choice上市公司财务数据合成的金融推理数据集，通过三阶段递进训练架构提升模型在复杂财务问答任务中的综合表现，目标在8B参数规模下实现显著的推理准确率提升。通过构建领域增强型推理机制，预期在现金流预测、财务指标关联分析等核心场景实现技术突破。

\subsection{研究内容}

本研究的主要工作集中在三个维度：在数据层面构建基于步骤级奖励的数据合成机制；在模型层面探索领域导向的过程奖励模型，并基于PEFT方法微调开源大模型；在推理层面建立基于PRM评分的偏好数据构建与DPO策略优化机制，提升复杂财务问题的可解释求解能力。通过系统性的技术创新，着力解决金融领域大模型应用中的知识注入不充分、数值推理不可靠等关键问题，为智能金融财务分析系统的实用化发展提供理论支撑与方法论指导。

\section{本文组织结构}

本文共分为五章，各章节的逻辑关系与主要内容如下：

\textbf{第一章 绪论}：阐述课题的研究背景与意义，综述金融财务推理、过程奖励模型、参数高效微调等方向的国内外研究现状，明确研究目标与研究内容，并给出本文的组织结构。

\textbf{第二章 相关技术与理论基础}：系统介绍本研究涉及的核心技术，包括大语言模型的Transformer架构与Decoder-Only范式、提示工程与思维链推理增强技术、LoRA等参数高效微调方法，以及奖励模型与强化学习在模型对齐中的应用，为后续章节的方法设计与实验研究提供理论支撑。

\textbf{第三章 总体设计方案}：从研究需求出发，提出"数据合成—过程监督—偏好对齐"的三阶段递进训练架构。详细阐述数据来源与处理流程、模型结构与工作机制、提示模板体系设计，以及步骤级与轨迹级双层奖励模型的联合训练机制。

\textbf{第四章 基于财务合成数据的过程奖励模型}：具体实施第三章的设计方案。在数据层面，基于Wind/Choice数据库通过Proposer→Solver→Validator三Agent协作管道合成金融推理数据集，并进行步骤级与轨迹级评分标注。在模型层面，构建Token分类式PRM架构，采用LoRA注意力微调策略（rank=16，仅训练0.75\%参数）。在实验层面，系统验证PRM训练收敛性、步骤级与轨迹级评估性能、置信度校准能力、过拟合修复效果，以及DPO偏好数据构建与策略模型训练的实验结果。

\textbf{第五章 结论与展望}：总结全文研究工作与核心创新点，客观分析当前方法的局限性，并展望未来的研究方向。

\newpage
\chapter{相关技术与理论基础}
\section{大语言模型基础}

大语言模型（Large Language Model，LLM）通常基于 Transformer 架构构建，其核心优势在于并行化的自注意力机制（Self-Attention），能够在大规模语料上学习到丰富的语义与结构信息。本节将介绍 Transformer 的基本结构，包括 Encoder–Decoder 框架、自注意力与多头注意力机制，以及近年来主流模型所采用的 Decoder-Only 架构。

\subsection{Transformer Encoder–Decoder 架构}

Transformer 最初由 Vaswani 等人在 2017 年提出\upcite{vaswaniAttentionAllYou2023}，其整体结构分为 Encoder 和 Decoder 两部分，每部分都由若干个堆叠的结构化模块组成。Encoder 负责对输入序列进行表示学习，Decoder 则在自回归模式下逐步生成目标序列。Transformer 整体架构如图~\ref{fig:transformer-arch} 所示。

Transformer Encoder 由多个相同结构的层堆叠构成，每层包含两个子模块：

1. 多头自注意力机制（Multi-Head Self-Attention）  
2. 前馈神经网络（Feed-Forward Network, FFN）

其核心思想是对输入序列 $\mathbf{X} = (x_1, x_2, ..., x_n)$ 进行并行自注意力建模，使模型能够基于整个上下文进行信息交互。

\subsection{Decoder}

Decoder 模块除包含与 Encoder 相同的自注意力与前馈结构外，还额外包含一个 Encoder–Decoder Attention，用于将已编码的输入序列信息融合到生成过程中。为了保证自回归生成，Decoder 内部的 Self-Attention 使用了 mask，确保第 $t$ 步仅依赖于之前的 $1..t-1$ 个 token。

\begin{figure}[H]
    \centering
    \includegraphics[width=0.8\textwidth]{Transformer.png}
    \caption{Transformer Encoder–Decoder 整体架构示意图}
    \label{fig:transformer-arch}
\end{figure}

\subsection{Encoder}

\subsubsection{自注意力机制与多头注意力机制}

Transformer 的核心是注意力机制（Attention）。给定查询（Query）、键（Key）和值（Value）矩阵 $Q,K,V$，Scaled Dot-Product Attention 的计算公式为：

\begin{equation}
\text{Attention}(Q, K, V) = \text{softmax}\left(\frac{QK^{T}}{\sqrt{d_k}}\right) V,
\end{equation}

其中 $d_k$ 为 Key 的维度，用于缩放内积防止数值过大。

\subsection{从自注意力到多头注意力}

自注意力（Self-Attention）指的是 $Q,K,V$ 都来自同一输入序列，从而实现序列内部所有 token 的信息交互。为了增强模型学习能力，Transformer 引入了多头注意力机制（Multi-Head Attention），其核心思想是将输入特征分成多个子空间，并行执行多组注意力计算，再将结果拼接。

多头注意力的数学表达式如下：

\begin{equation}
\text{MultiHead}(Q,K,V)=\text{Concat}(head_1, ..., head_h) W^O,
\end{equation}

其中每个 $head_i$ 为：

\begin{equation}
head_i = \text{Attention}(QW_i^Q, KW_i^K, VW_i^V),
\end{equation}

$W_i^Q, W_i^K, W_i^V, W^O$ 均为可学习参数。

图\ref{fig:multihead-attn} 展示了多头注意力机制结构。

\begin{figure}[H]
    \centering
    \includegraphics[height=0.3\textwidth]{MultiHead.png}
    \caption{多头注意力机制示意图}
    \label{fig:multihead-attn}
\end{figure}

多头机制允许模型在不同子空间学习不同的语义与结构关系，从而显著提升表达能力。

\subsection{Decoder-Only 架构}

随着 GPT 系列模型的发展，主流大语言模型逐渐转向了仅使用 Transformer Decoder 堆叠的 Decoder-Only 架构。这类架构省去了 Encoder 模块，全部输入通过自回归方式处理，并严格使用 Masked Self-Attention。

给定输入序列 $x_1, ..., x_T$，模型在第 $t$ 步仅利用先前 token 的表示：

\begin{equation}
P(x_t|x_{<t}) = \text{TransformerDecoder}(x_{1:t-1})
\end{equation}

最终目标是最大化序列生成的似然：

\begin{equation}
\mathcal{L} = -\sum_{t=1}^{T} \log P(x_t | x_{<t})
\end{equation}

该结构具有以下优点：

1. 训练与推理过程一致（均为自回归生成）
2. 无需 Encoder，计算图更简单
3. 适合大规模预训练与扩展
4. 推理更高效，适合实际应用

GPT 系列模型率先采用此架构：GPT-1 通过生成式预训练实现多任务迁移\upcite{radfordImprovingLanguageUnderstanding2018}；GPT-3 进一步验证了大规模 Decoder-Only 模型在少样本学习中的强大性能\upcite{brownLanguageModelsAre2020a}。随后，LLaMA\upcite{touvronLLaMAOpenEfficient2023}、Qwen\upcite{baiQwenTechnicalReport2023}、DeepSeek 等开源大模型均采用类似的纯 Decoder 架构，使其成为当前 LLM 的事实标准。
\section{提示工程（Prompt Engineering）}

提示工程（Prompt Engineering）是指通过设计合理的输入提示（Prompt），引导大语言模型（LLM）生成符合预期的输出内容的技术。在大语言模型中，Prompt 的设计直接影响生成质量、内容准确性以及推理能力。本节将从模型采样参数、零样本/单样本/少样本提示、指令调优、基础提示技巧以及思维链等方面进行说明。

\subsection{模型采样参数}

在生成文本时，模型通常会基于概率分布选择下一个 token。假设在时间步 $t$，模型对词汇表中每个 token 的预测概率为 $P(x_t|x_{<t})$。为了控制生成文本的多样性与创造性，引入以下采样参数：

\subsubsection{Temperature}

Temperature 参数 $\tau$ 调节概率分布的平滑程度：

\begin{equation}
P_\text{temp}(x_t|x_{<t}) = \frac{P(x_t|x_{<t})^{1/\tau}}{\sum_{x'} P(x'|x_{<t})^{1/\tau}}
\end{equation}

- $\tau < 1$：概率分布尖锐，更倾向高概率 token（保守生成）  
- $\tau > 1$：概率分布平滑，更倾向多样化生成

\subsubsection{Top-k 采样}

Top-k 采样限制候选 token 为概率最高的 $k$ 个：

\begin{equation}
P_\text{top-k}(x_t|x_{<t}) =
\begin{cases}
\frac{P(x_t|x_{<t})}{\sum_{i \in \text{Top-}k} P(x_i|x_{<t})}, & x_t \in \text{Top-}k \\
0, & \text{otherwise}
\end{cases}
\end{equation}

\subsubsection{Top-p (Nucleus) 采样}

Top-p 采样选择概率累计超过阈值 $p$ 的最小集合 $\mathcal{V}_p$：

\begin{equation}
\mathcal{V}_p = \min \left\{ V' \subseteq V : \sum_{x_i \in V'} P(x_i|x_{<t}) \ge p \right\}
\end{equation}

然后在 $\mathcal{V}_p$ 中进行归一化采样。该方法被证明能有效缓解神经文本退化问题，生成更具多样性且连贯的文本\upcite{holtzmanCuriousCaseNeural2020}。


\subsection{零样本、单样本与少样本提示}

提示工程根据是否提供示例，可分为三类：

1. 零样本（Zero-Shot）：仅提供任务描述，无任何示例  
   例：\texttt{“请将下面句子翻译成英文。”}

2. 单样本（One-Shot）：提供一个输入-输出示例作为参考  
   例：\texttt{“例句：‘你好’ -> ‘Hello’。请翻译：‘再见’。”}

3. 少样本（Few-Shot）：提供多个示例，增强模型理解任务的能力  
   少样本提示通常显著提升复杂任务的表现，尤其是逻辑推理或代码生成任务\upcite{weiFinetunedLanguageModels2022}。


\subsection{指令调优的影响}

指令调优（Instruction Tuning）是指在大规模模型上，使用明确任务指令和示例数据进行额外训练，使模型更好地理解和执行自然语言指令。经过指令调优的模型在零样本和少样本条件下表现更稳定，更容易被 Prompt 引导生成准确结果\upcite{weiFinetunedLanguageModels2022}。


\subsection{基础提示技巧}

1. 明确任务描述：使用清晰、具体的自然语言描述任务  
2. 提供结构化格式：如表格、编号或 JSON 格式，便于模型输出可解析内容  
3. 指定风格或长度：如“请用学术风格回答，不超过 100 字”  
4. 分步提示：将复杂任务拆分为子任务，逐步生成结果


\subsection{思维链（Chain-of-Thought, CoT）}

思维链提示是一种通过引导模型显式生成推理过程来提升复杂任务表现的方法。其核心思想是：

\begin{equation}
x_\text{input} \longrightarrow \text{推理步骤}_1, \text{推理步骤}_2, ..., \text{最终答案}
\end{equation}

如：\texttt{“请先写出解题步骤，再给出最终答案。”}

通过显式生成中间推理步骤，模型在数学推理、逻辑推理等任务中的准确率显著提升\upcite{weiChainofThoughtPromptingElicits2023}。
\section{推理增强技术}

推理增强技术旨在提升大语言模型在复杂任务中的推理能力，使其能够生成更加准确和有逻辑性的输出。传统的大语言模型在处理复杂问题时，直接生成最终答案往往容易出现错误或者遗漏关键推理步骤。为了解决这一问题，研究者提出了多种增强推理的方法，其中包括思维链（Chain-of-Thought, CoT）、ReAct 模型以及基于检索增强的生成方法（Retrieval-Augmented Generation, RAG）。

思维链技术通过引导模型显式地生成中间推理步骤，使模型在推理过程中形成类似人类的思考路径。例如，在数学推理或逻辑题中，模型不仅输出最终答案，还需要逐步推导每一步的结果。形式上，如果输入为 \(x\)，模型生成的输出序列可表示为：
\begin{equation}
y = \text{Step}_1 \rightarrow \text{Step}_2 \rightarrow \dots \rightarrow \text{Step}_n \rightarrow \text{Answer}
\end{equation}
通过显式的推理步骤，模型在多步逻辑、条件判断和算术推理任务中的性能显著提高。此外，思维链提示也为后续复杂任务的自我验证（self-verification）和迭代推理提供了基础，使模型能够在生成答案时自行发现潜在错误并进行修正。

ReAct 模型将推理（Reasoning）与行动（Acting）结合，通过在推理过程中引入外部环境交互或工具调用，实现更高效的任务完成能力。在该框架下，模型的生成过程不仅包含内部推理步骤，还能够根据环境反馈进行动态调整。例如，模型在回答一个包含知识库查询的问题时，先生成推理步骤，再调用检索工具获取必要信息，最后整合生成答案。形式上，可表示为：
\begin{equation}
y_t = f_\theta(x_t, s_{t-1}, e_{t-1})
\end{equation}
其中 \(s_{t-1}\) 表示模型的内部状态，\(e_{t-1}\) 表示环境反馈，函数 \(f_\theta\) 表示模型的生成函数。通过引入环境交互，ReAct 模型在问答、规划和决策任务中表现出更强的适应性和准确性。

检索增强生成（RAG）方法通过结合外部知识库，使模型在生成过程中能够访问大规模事实信息，从而提升事实性和专业性输出。与纯生成式模型相比，RAG 先根据输入从知识库中检索相关文档，然后将检索到的内容作为额外上下文提供给语言模型生成答案。其生成过程可表示为：
\begin{equation}
y \sim P_\theta(y|x, \mathcal{R}(x))
\end{equation}
其中 \(\mathcal{R}(x)\) 表示从知识库检索到与输入 \(x\) 相关的内容。RAG 方法在金融、医疗和科研问答等需要高事实准确性的场景中尤为有效，能够显著减少虚构信息的产生。

\section{PEFT 微调方法}

随着大语言模型规模的不断扩大，传统的全参数微调方法面临着显著的计算和存储成本问题。在这种背景下，参数高效微调（Parameter-Efficient Fine-Tuning, PEFT）方法逐渐成为主流。PEFT 方法的核心思想是，在保留预训练模型大部分参数不变的情况下，仅对少量参数进行调整，从而显著降低微调的计算开销，同时仍能在下游任务中取得良好效果。其中，LoRA（Low-Rank Adaptation）是一种代表性的 PEFT 技术，其在理论基础与实际应用中都展现出了极高的效率和灵活性。

LoRA 的基本原理是对预训练模型的权重矩阵进行低秩分解，从而只学习少量附加参数，而保持原始权重冻结不变。假设某一线性层的权重矩阵为 \(W_0 \in \mathbb{R}^{d_\text{out} \times d_\text{in}}\)，LoRA 将权重更新表示为两个低秩矩阵的乘积：
\begin{equation}
\Delta W = A B, \quad A \in \mathbb{R}^{d_\text{out} \times r}, \quad B \in \mathbb{R}^{r \times d_\text{in}}
\end{equation}
其中 \(r \ll \min(d_\text{out}, d_\text{in})\) 是低秩维度，通常远小于原权重矩阵的维度。在微调过程中，模型实际使用的权重为：
\begin{equation}
W = W_0 + \alpha \Delta W
\end{equation}
其中 \(\alpha\) 是缩放系数，用于控制低秩更新对整体权重的影响。通过这种方式，LoRA 仅需训练 \(r(d_\text{in} + d_\text{out})\) 个参数，而不需要更新原始的 \(d_\text{in} \times d_\text{out}\) 个参数，从而大幅度降低存储与计算成本。

LoRA 的优势不仅在于参数量的减少，还体现在与预训练模型结构的兼容性。它可以直接嵌入到 Transformer 的自注意力层或者前馈网络层中，尤其适用于大模型的微调。对于自注意力层中的查询（Q）、键（K）、值（V）矩阵，LoRA 通常仅对查询和键矩阵进行低秩适配，而保持值矩阵不变，这样既保证了生成性能，又避免了过度拟合。此外，LoRA 可以灵活组合其他 PEFT 方法，如适配器（Adapter）和前馈层微调，使得微调策略更加多样化。下图展示了 LoRA 的整体架构：

\begin{figure}[h!]
\centering
\includegraphics[width=0.8\textwidth]{LoRA.PNG}
\caption{LoRA 低秩适配架构示意图}
\label{fig:lora}
\end{figure}

在训练目标方面，LoRA 通常采用标准的监督学习或对比学习损失，但只对低秩矩阵 \(A, B\) 进行梯度更新。假设下游任务的损失函数为 \(\mathcal{L}(\theta)\)，则 LoRA 微调的优化问题可形式化为：
\begin{equation}
\min_{A,B} \mathcal{L}(W_0 + \alpha AB)
\end{equation}
这种结构使得优化空间大幅减小，同时保留了预训练模型丰富的语言知识，从而在少量训练数据下仍能获得良好的下游性能。

除了 LoRA，PEFT 领域还出现了多种方法，例如 Adapter、Prefix Tuning 和 Prompt Tuning。Adapter 在每一层插入小型瓶颈网络，对原始隐藏层进行线性变换后再加回原输出，从而实现微调；Prefix Tuning 则通过训练一组可学习的前缀向量，作为 Transformer 的额外输入，引导模型生成符合任务需求的输出；Prompt Tuning 更侧重于优化输入提示（prompt），在冻结模型参数的前提下调整可训练的嵌入向量，以适应下游任务。这些方法与 LoRA 共享的核心理念是参数高效、低成本，同时保持预训练模型的语言理解能力。

在实际应用中，LoRA 的低秩适配方法特别适合大规模语言模型的多任务学习和跨领域迁移。例如，在金融、医疗等专业领域，数据标注成本高且数量有限，LoRA 能够在有限样本条件下实现模型微调，同时保持模型生成的准确性和可靠性。此外，LoRA 与 RLHF 等训练方法兼容，可以结合奖励信号进一步优化模型输出，使其在特定任务中更加符合人类偏好。由于其灵活性和高效性，LoRA 已成为当前 PEFT 微调研究和实践中的首选方案。

总的来看，PEFT 微调方法通过减少可训练参数、降低计算负担并保持预训练模型知识的完整性，为大语言模型的高效微调提供了理论和实践支持。其中 LoRA 以低秩分解为核心，通过显著减少参数量实现微调目标，兼顾性能与效率，已成为学术界和工业界广泛采用的微调策略。这一方法不仅推动了大模型在资源受限环境下的应用，也为未来更复杂任务的高效微调提供了坚实基础。


\section{奖励模型与强化学习}

在大语言模型的优化过程中，单纯的最大似然训练往往难以保证模型输出符合人类偏好或特定任务目标。为此，引入奖励模型（Preference Reward Model, PRM）与强化学习（Reinforcement Learning, RL）成为提高模型对齐性能的重要手段。奖励模型通过学习人类对输出的偏好，对模型生成结果进行评分，从而为强化学习提供优化目标。

PRM 的核心在于构建一个可训练的评分函数 \(R_\phi(y, x)\)，用于衡量模型在给定输入 \(x\) 下生成输出 \(y\) 的质量。通常通过人类标注的数据集进行训练，优化目标为最大化人类偏好对的概率：
\begin{equation}
\mathcal{L}_{PRM}(\phi) = - \sum_{(y_+, y_-, x) \in D} \log \sigma(R_\phi(y_+, x) - R_\phi(y_-, x))
\end{equation}
其中 \(y_+\) 与 \(y_-\) 分别表示优选输出与次优输出，\(\sigma\) 为 Sigmoid 函数，确保奖励模型能够区分更优的答案。

在训练完成奖励模型后，强化学习被用来进一步优化生成模型，使其输出更加符合人类偏好。典型方法包括 Proximal Policy Optimization (PPO) 和 Direct Preference Optimization (DPO)\upcite{rafailovDirectPreferenceOptimization2024}。在 PPO 框架下，模型被视为策略 \(\pi_\theta(y|x)\)，通过最大化带有奖励的期望回报进行更新：
\begin{equation}
\mathcal{L}_{PPO}(\theta) = \mathbb{E}_{y \sim \pi_\theta} \left[ \frac{\pi_\theta(y|x)}{\pi_{\theta_\text{old}}(y|x)} \hat{A}(y,x) \right]
\end{equation}
其中 \(\hat{A}(y,x)\) 表示优势函数，衡量当前生成相比于基准策略的改善程度。PPO 通过限制策略更新幅度，保证训练过程稳定且收敛可靠。DPO 则通过直接将奖励模型信号融入损失函数，优化模型生成与偏好分布的一致性，其核心思想是最大化生成输出在奖励模型下的偏好概率。

除了 PPO 和 DPO，RLHF（Reinforcement Learning from Human Feedback）通常结合预训练语言模型、奖励模型和策略优化，形成完整训练流程。该方法首先利用人类标注数据训练奖励模型，然后用强化学习优化生成策略，使输出在符合语法、逻辑的同时，也最大化人类偏好评分。形式化地，RLHF 的目标可表示为：
\begin{equation}
\max_\theta \mathbb{E}_{y \sim \pi_\theta} \big[ R_\phi(y, x) - \beta \cdot \text{KL}(\pi_\theta(y|x) || \pi_\text{ref}(y|x)) \big]
\end{equation}
其中 \(\beta\) 为权衡因子，\(\pi_\text{ref}\) 为预训练基准模型，KL 项用于防止策略偏离原始语言分布过远，保证生成语言的自然性和多样性。

通过奖励模型提供的偏好信号与强化学习策略优化，PRM 与 RLHF 方法有效地提升了大语言模型的对齐能力，使模型在生成高质量、符合人类期望的文本方面表现显著优于传统训练方法。

\section{本章小结}

本章系统介绍了大语言模型及其相关技术基础，为后续章节的研究提供了理论支撑。首先，我们从 Transformer 架构出发，详细阐述了 Encoder–Decoder 与 Decoder-Only 两种主要结构，以及自注意力和多头注意力机制的原理与公式。这些机制奠定了大语言模型强大表示能力和并行计算能力的基础，使模型能够在大规模语料上高效学习复杂的语义与结构关系。

随后，针对模型的实际使用场景，介绍了提示工程技术。通过对采样参数的调节、零样本、单样本与少样本提示的应用，以及指令调优与基础提示技巧的说明，本章展示了如何通过合理设计输入引导模型生成符合预期的输出。特别是思维链（Chain-of-Thought）技术的引入，使模型能够在推理任务中显式生成中间步骤，从而显著提升复杂任务的准确性和可靠性。

在推理增强技术方面，本章探讨了思维链、ReAct 模型和检索增强生成（RAG）等方法。这些方法不仅提升了模型的逻辑推理能力，也增强了模型在面对需要外部知识和多步决策任务时的适应性和准确性，为后续实际应用提供了坚实基础。

最后，本章介绍了奖励模型与强化学习（PRM 与 RLHF）在模型对齐中的作用。通过建立人类偏好奖励模型并结合强化学习策略优化，模型能够在保持语言生成质量的同时，更好地符合人类预期和任务要求。这一技术框架为大语言模型在对话系统、专业问答及决策支持等领域的应用提供了关键支持。

综上所述，本章不仅梳理了大语言模型的核心结构与原理，也系统介绍了引导生成、增强推理及优化对齐的关键技术，为后续章节对模型应用与改进的深入研究奠定了理论与方法基础。


\newpage
\chapter{总体设计方案}
\section{研究需求}
本课题面向银行金融业务的智能化财务分析需求，旨在构建具备领域适应性和可解释性的金融推理模型体系。现有大语言模型在财务文本理解、数值计算、跨模态信息融合以及合规性生成方面仍存在显著不足，难以支撑高风险、强逻辑的金融决策场景。因此，本研究的总体需求集中体现在以下三个方向。

首先，需要构建覆盖金融知识结构、推理步骤链路与真实业务情境的数据体系。金融任务具有知识密集、概念关联复杂和推理链条较长的特点，单纯依赖通用语料难以支撑高质量训练。模型必须具备从专家分析中提取金融知识、理解会计科目关系、处理多步骤计算逻辑的能力。此外，基于过程级奖励的数据需求强烈，需要能够覆盖“正确步骤、错误步骤、部分正确步骤”的多类型样本，以支持细粒度监督。

其次，需要形成适配金融领域的推理机制与模型训练方法。现有大模型在财务推理中容易出现数值错误、逻辑跳跃与事实偏差，特别是在现金流预测、财务指标关联分析等场景下。为提升可靠性，需要引入领域导向的奖励模型，对推理链路的逻辑性、知识使用情况和计算正确性进行评估。同时，引入SFT监督微调和强化学习体系，降低专用模型的训练成本，实现小规模参数的高效适配。在推理层面，需要配备具备解释能力的 Agent 机制，并结合 RAG 检索强化模型的金融知识使用能力。

最后，需要建立可验证、可对齐的任务求解流程。金融决策场景对生成内容的可信度和合规性要求极高，因此模型必须具备可验证推理路径。需求包括：能够追踪每一步推理行为；能够对答案来源、知识引用以及计算逻辑进行回溯；能够通过奖励信号实时反馈推理质量。为此，需要设计标准化提示模板、分阶段训练公式以及奖励模型的迭代策略，确保模型在复杂任务下保持稳定性与一致性。

综合来看，研究需求指向“数据—模型—推理”三个层级的系统构建，目标是在通用大模型基础上形成适用于金融领域的高可靠推理体系，支撑财务分析、合规判断与管理决策等关键业务。
\section{总体研究路线}
如图所示~\ref{fig:pipeline}，本研究总体设计方案涵盖数据构建、奖励模型训练与推理优化三大核心模块，形成闭环的金融推理模型开发流程。我们从基座模型开始，冻结其全部预训练参数，使其语义能力保持不变。在此基础上添加 Reward Head，让模型可以在每个步骤位置输出奖励分数。为了适应金融推理任务，我们通过 LoRA 修改注意力层，使模型的注意力模式逐步偏向“正确推理步骤”“关键概念出现点”等位置。训练过程使用两个损失项：步骤级损失和轨迹级损失，使模型既能判断局部逻辑是否合理，又能判断整体推理方向是否正确。最终得到的 PRM 能够对任何推理轨迹进行评分，为后续策略优化提供核心评价标准。

本研究整体采用\textbf{三阶段递进训练架构}：（1）阶段一，训练基于Token分类的PRM过程奖励模型，实现步骤级过程监督（步骤准确率达95\%）；（2）阶段二，基于PRM对候选推理轨迹进行评分排序，构建DPO偏好数据集（平均margin 0.686）；（3）阶段三，采用DPO偏好对齐训练优化金融推理策略模型。该架构将传统GRPO在线奖励搜索范式替换为DPO离线偏好对齐，训练流程更简洁稳定。

\begin{figure}[H]
    \centering
    \includegraphics[height=0.8\textwidth]{pic/pipeline.png}
    \caption{模型整体架构示意图}
    \label{fig:pipeline}
\end{figure}


\subsection{数据来源与处理流程}

本研究的数据来源与处理流程围绕金融领域财务推理任务展开，旨在为过程奖励模型训练构建高质量、可监督的推理数据集。整体流程分为两个阶段：第一阶段为QA问答对合成，第二阶段为评分标注。

\subsubsection{数据来源}

本研究以Wind/Choice上市公司财务数据库作为主要数据来源，提取每条记录的财务快照，包含核心指标、补充指标、原始财务字段、历史指标和趋势特征五类指标。数据覆盖盈利能力、偿债能力、营运能力、结构性风险和现金流质量五个评估维度，为后续推理轨迹生成与奖励标注提供结构化的领域知识基础。

\subsubsection{两阶段数据合成概述}

第一阶段为QA问答对合成：基于Proposer→Solver→Validator三Agent协作管道，在LangGraph状态图上循环迭代，自动生成包含正负样本的推理轨迹数据。第二阶段为评分标注：对每条推理轨迹进行步骤级与轨迹级两级评分，通过LLM评判生成0/1二值标签，构建可监督的过程奖励数据集。

两阶段的详细设计与实施参见第4章。


\subsection{模型结构与工作机制}

本研究的模型结构围绕金融推理场景展开，目标是让模型能够识别推理轨迹中的高价值步骤，并理解轨迹整体的质量。模型以预训练语言模型为基础，并通过过程奖励建模的方式进行改造。整体结构强调可解释性、可控性与金融领域的知识一致性。模型训练强调轨迹结构化表达、标签精确对齐与推理片段的显式监督，以提升模型的推理评估能力。

模型结构由基础语言模型、奖励头、特殊符号机制以及数据对齐组件构成。工作机制包括推理轨迹序列构造、标注信号对齐、训练阶段优化以及分布式训练等环节。整体框架强调模块化设计，使模型能够处理长度较长、结构复杂的金融推理数据。

\subsubsection{基础模型结构与参数初始化}

基础模型采用通用语言模型作为底座。本研究选取规模适中的模型，并在其上构建过程奖励模块。奖励模块基于 Transformer 框架实现，以更好适应预训练模型的表示方式。模型加载基础参数后，再根据任务需求引入新的结构。该方式既保持通用语义能力，又为后续金融推理场景提供优化空间。

模型使用独立配置文件定义结构，包括隐藏层大小、注意力头数量以及奖励头维度等参数。这种设计使模型更易扩展，也便于在不同规模间迁移。模型加载预训练权重后，通过严格对齐方式将参数映射到奖励模型内部结构，确保模型能够在新任务上稳定启动。

\subsubsection{过程奖励头的设计思路}

本研究关注的是推理过程的质量，而不是单一结果。因此模型在原始语言建模结构外增加奖励头，用于输出步骤级奖励与轨迹级奖励。奖励头为一个轻量化2层MLP（Linear→ReLU→Linear→2分类logits），读取Transformer在特殊符号位置的隐藏状态，映射为二分类质量评分（0/1）。训练时仅在这些特殊符号位置计算交叉熵损失，其余位置忽略。

奖励头结构保持轻量化，以减少额外计算开销，同时保持了对推理片段的敏感性，提供更强的过程监督能力。模型输出的每个奖励值与训练数据中的质量标签严格对齐，以构成完整的监督链路。

\subsubsection{特殊符号标记机制}

推理轨迹具有清晰的结构。本研究引入两个特殊符号来表达这些结构：一个表示步骤分界，另一个表示轨迹结束。该机制使模型能够区分自然语言内容与推理结构内容，使轨迹可被模型显式理解。每个分界符对应一个标签位置，从而建立步骤级奖励的监督信号。

这种设计让模型具备识别推理边界的能力，使其可以根据真实结构进行判断。推理轨迹不再被视为普通文本，而是具有分块特征的结构化序列。该特性对金融推理任务尤其重要，因为许多推导过程依赖连贯的中间步骤，而不是单一最终结果。

\subsubsection{数据对齐机制与输入构造}

为了让模型正确接收奖励监督，本研究设计了输入数据构造方式。每条样本包含问题、选项、推理轨迹以及最终答案。这些内容按照固定模板组织成为统一序列，以保证模型输入结构稳定。推理轨迹中的换行被替换为特殊符号，使模型能够感知步骤边界。

数据对齐过程中，标签序列与输入序列严格对应。步骤标签被放置在分界符位置，轨迹标签放置在最终结束符位置。其余位置填充无效标签，从而避免模型误学习非目标位置。该机制保证了监督信号的稀疏性与精准性，使训练更加稳定。

\subsubsection{金融推理场景下的表达能力增强}

在金融领域中，问题往往涉及数据计算、概念解释、逻辑推导以及业务规则。模型在训练过程中通过奖励机制学习如何识别合理路径，因此在结构上必须支持长序列输入。本研究将最大长度设置为较大值，使模型能够处理完整的推理轨迹，而非只处理其片段。

模型在训练中保持对原有语义建模能力的继承，在此基础上进一步强化金融推理表达能力。模型不仅识别语言，也识别推理步骤间的逻辑关系，从而具备更强的结构理解能力。这种设计使模型能够从长推理链中提取关键信息并判断其质量。

\subsubsection{训练工作机制}

模型训练采用监督微调方式，通过奖励标签引导模型学习奖励函数。训练框架基于专用训练器实现，支持长序列、稀疏标签以及结构化文本输入。训练器根据任务需求选择合适的优化器与调度策略，以提升学习质量。

训练过程中，数据按照批次输入，模型输出对应的奖励值。损失函数只计算特殊符号位置的误差，以突出关键位置的学习目标。训练循环逐步调整模型参数，使其能够更精准地识别推理中的高质量步骤。模型在多个周期内反复优化，使性能逐渐收敛。

\subsubsection{训练资源与优化策略}

训练在单张NVIDIA RTX PRO 6000 Blackwell Server Edition GPU（98GB显存）上进行，大显存设计消除了跨设备通信开销。训练过程中采用梯度累积与bf16混合精度以提升训练效率，配合8bit AdamW优化器进一步降低显存占用，使模型能够在单卡上稳定完成长序列金融推理数据的训练。

\subsubsection{整体工作机制总结}

模型结构与工作机制形成一个完整链条。模型首先通过结构化输入理解推理轨迹的组成，再通过奖励头对关键位置进行评分。训练阶段通过精确标签对齐引导模型学习质量评估能力。该训练机制使上述流程能够高效地扩展到更大规模的训练场景。

整体机制强调三个核心目标：保持预训练模型的基本能力、引入金融推理过程的显式监督、实现结构化推理片段的识别能力。这些目标通过统一的结构设计得以实现，使模型能够在金融任务中执行更稳健、更可解释的推理质量判断。

\subsection{提示模板设计}

本研究的任务核心在于构建一个能够理解金融场景、识别推理步骤质量、并完成长链条推理评估的过程奖励模型。因此，提示模板的设计不仅服务于生成推理轨迹，还承担着引导模型完成质量判断、知识抽取与事实核验等多重任务。提示模板的结构直接影响数据质量，进而影响奖励模型的学习效果。本节对提示模板的整体设计思路进行概述，强调其逻辑结构、功能定位与信息组织方式。

提示模板体系由多类指令构成，每类模板负责不同子任务。其设计思路遵循三个原则：一是保持任务边界清晰，使模型能够理解角色与责任；二是强调结构化输出，使标签对齐更加稳定；三是引导模型进行显式推理，使推理链具备可监督性。模板不会依赖外部知识，而是围绕问题本身组织推理过程。以下从整体架构、模块关系与生成逻辑三个角度进行说明。

\subsubsection{整体提示体系的结构化设计思路}

整个提示体系围绕推理任务展开。从数据生成到质量评价，再到知识抽取和事实核验，均需要专门的提示模板来引导模型输出结构化内容。模板整体结构呈现多层次特征：

首先，模板明确模型的角色定位。例如，有些模板要求模型扮演推理专家，有些要求扮演逻辑检验者，还有些要求扮演金融知识抽取者。这种角色分工使模型的思考模式更贴近实际任务。

其次，模板以高度结构化的格式组织输入信息。无论是推理轨迹、专家分析还是知识点，均按照固定格式呈现。这样可以保证模型在不同样本上保持一致的解析方式。

最后，模板使用明确的信号标记，包括思想部分、答案部分以及 JSON 格式输出，使得输出结果可以被自动解析、提取并用于训练。通过这种方式，模板成为连接数据生成与模型学习的关键纽带。

\subsubsection{思维链生成模板的设计理念}

推理轨迹生成模板是整个系统的起点。此模板引导模型从问题出发，构建完整的推理链。设计此模板的目的不是让模型直接给出答案，而是让其暴露推理路径，从而提供可监督的轨迹。

模板中使用了“Thought”与“Solution”的双结构，使模型先展示完整推理过程，再给出最终答案。这种结构通过显式分层，将推理与答案隔离开来，使两者可以分别使用不同的监督方式。模型在“Thought”部分展开细致思考，其中每一步都采用统一格式标记。此设计使得每个步骤可独立评价，也便于与奖励模型的标注信号对齐。

该模板还要求模型进行反思式推理。其设计理念是让模型在推理中修正前序步骤的不足，使轨迹更加接近真实的逻辑推导过程。此方式增强了数据多样性，使过程奖励模型能够学习更丰富的推理结构。

\subsubsection{金融知识抽取模板的设计思路}

知识抽取模板的作用在于从专家分析中提取关键概念，并为每个概念生成定义。其设计强调两点：一是识别关键术语，二是用简洁方式组织解释。

模板要求输出为 JSON 列表，以保证后续处理的自动化和标准化。每个知识点包含术语与解释两个部分。这样的结构符合知识库的组织方式，也使得后续事实校验任务能够直接引用这些知识点。

模板以专家分析为输入，从中提取结构化概念。设计中使用自然语言指令引导模型关注金融核心概念，而非文本细节，使模型的输出更聚焦于知识本身。

\subsubsection{步骤质量评价模板的设计理念}

质量评价模板由三个子模板组成，分别对应逻辑性、步骤准确性与目标推进性等维度。该设计遵循多角度评价原则，使模型能够从不同视角判断步骤质量。

\begin{itemize}
    \item 逻辑评价模板引导模型关注步骤的逻辑结构，包括是否存在推理跳跃或逻辑断裂。
    \item 步骤正确性模板要求模型评估内容本身是否成立，从而排除错误事实或不当推导。
    \item 目标推进性模板则要求模型判断该步骤是否朝最终答案方向推进，以衡量步骤在整体任务中的作用。
\end{itemize}

所有评价模板均使用 JSON 输出，使每个维度得到明确量化。这对于步骤级奖励的构建至关重要。

\subsubsection{程序正确性模板的设计原因}

程序正确性模板专注于判断某一步是否向正确答案推进。此模板的设计逻辑强调“方向性”，而不是逻辑本身。其作用在于让模型理解一个关键事实：即便某步骤逻辑合理，但若与最终答案方向不一致，也不应被视为有效步骤。

模板输入包括正确答案，使模型能够进行目标对照。输出采用二值标注，使得标签更加稳定。该模板为奖励模型提供了非常重要的监督信号，有助于让模型理解步骤的重要性来源于其对最终目标的贡献。

\subsubsection{事实核验模板的设计理念}

事实核验模板围绕知识库展开。其核心任务是检查当前步骤是否与金融知识一致。模板提供知识库内容，让模型作为事实验证者进行判断。

此模板的设计强调严格性与一致性。模型需要比对步骤中的术语与知识库中的术语，并判断是否匹配。这一过程使得模型能够区分“事实性正确”与“逻辑性正确”。对于金融推理任务而言，这一区分至关重要。

模板输出为二值结构，使得奖励模型训练更加直接。该设计强化了模型识别金融错误的能力，有助于提高整体推理可靠性。

\subsubsection{答案比对模板的功能定位}

答案比对模板的作用是确认模型生成的最终答案是否正确。模板采用结构化方式对答案进行比对，使得轨迹级奖励能够准确定义。

其设计不关注推理过程本身，而只关注最终答案的匹配性。这种设计使轨迹层面的评分不受表述差异影响，从而保证了评价的稳定性。

\subsubsection{知识覆盖度模板的设计价值}

知识覆盖度模板用于衡量模型在推理轨迹中是否使用了应使用的知识点。该模板连接了知识抽取与推理轨迹，为轨迹级奖励提供了知识基础。

模板要求模型输出覆盖数量与对应索引，使得每条轨迹可以量化评价其知识利用程度。这种设计强化了模型对知识的调用能力，使推理轨迹更贴近真实金融业务分析过程。

\subsubsection{提示模板体系的整体协同机制}

整个提示体系协同运作，构成完整的推理数据生成与质量评价闭环。推理生成模板负责提供轨迹数据。质量评价模板负责提供步骤级信号。答案匹配与知识覆盖模板提供轨迹级信号。它们共同构成奖励模型训练所需的多维监督框架。

这种提示体系的设计使得模型不仅能够生成推理，还能够理解推理、评价推理并校验推理。模板之间相互依赖并形成逻辑链条，使金融推理任务从数据到模型均具备内在一致性。


\subsection{奖励模型的设计与迭代机制}

本节围绕奖励模型的构建方法与迭代机制展开说明。内容来源于前述训练框架，同时结合实际代码结构加以修正与完善。整体设计强调过程监督、结构化奖励构成与可重复的训练流程。模型在训练中以思维轨迹为核心对象，通过对关键标记位置进行监督，使其能对每一步推理行为给出判断。整个机制以“可解释”“可拆解”“可迭代”为基本设计逻辑。

奖励模型的任务目标是对推理过程进行细粒度评价。模型需要理解问题文本、推理轨迹以及最终答案，并在此基础上产出两类奖励：一类用于步骤级判断，另一类用于整体轨迹判断。前者反映局部质量，后者反映整体方向。两者共同构成奖励体系的核心结构，并在训练时被整合进联合损失函数。

\subsubsection{总体设计逻辑}

奖励模型的设计遵循三个原则。第一，奖励来源应可追溯。第二，不同层级的奖励应保持相对独立，以便模型学习局部信息与整体信息。第三，训练流程需要与数据构造方式一致，确保标签与位置对应关系明确。基于这三个原则，模型结构与奖励机制被组织为以下几个部分：

\begin{itemize}
    \item 步骤级奖励，用于判断每一个推理片段的质量；
    \item 轨迹级奖励，用于判断整体推理过程是否合理；
    \item 奖励聚合策略，用于整合不同来源的信号；
    \item 训练目标函数，用于约束模型同时学习两个层级的判断。
\end{itemize}

这些部分在代码中以数据构造函数、模型结构定义以及训练器配置的形式体现。尤其在数据整理阶段，程序会把奖励标签与推理片段的分隔符绑定。模型读取输入后，可以根据标记位置对输出进行监督，从而使奖励模型具备识别推理过程的能力。

\subsubsection{步骤级奖励的构成}

步骤级奖励反映推理过程中每一步的表现质量。其设计思路是把复杂的推理行为拆解成多个可观察片段，每个片段在数据中通过特殊标记进行定位。在此基础上，通过LLM对每一个片段进行三个维度的评分，最终整合为单一的步骤级二值标签。

步骤级奖励的三个评分维度分别为定性评分、过程正确性评分以及事实准确性评分。

\textbf{定性评分（\(r_{\mathrm{qual}}\)）}：LLM从三个子维度对步骤进行0--1连续评分——逻辑合理性（logical\_soundness）、步骤正确性（step\_correctness）、目标推进性（target\_progression），取三者均值：
\[
r_{\mathrm{qual}} = \frac{1}{3}(\text{logical\_soundness} + \text{step\_correctness} + \text{target\_progression})
\]

\textbf{过程正确性（\(r_{\mathrm{procedural}}\)）}：LLM接收问题、推理历史和已知正确答案，判断当前步骤是否为朝向正确答案的合理逻辑推进，输出二值标签：
\[
r_{\mathrm{procedural}} \in \{0, 1\}
\]

\textbf{事实准确性（\(r_{\mathrm{factual}}\)）}：LLM接收展平的financial\_data指标列表和当前步骤文本，验证步骤中使用的所有指标数值是否与知识库一致，输出二值标签：
\[
r_{\mathrm{factual}} \in \{0, 1\}
\]

首先合并过程正确性和事实准确性得到准确性得分：
\[
r_{\mathrm{acc}} = 0.5 \times r_{\mathrm{procedural}} + \omega_k \cdot r_{\mathrm{factual}}, \quad \omega_k = 1.0
\]

然后将定性评分与准确性得分通过Softmax动态加权聚合：
\[
r_t^{\mathrm{step}} = \mathrm{softmax}([r_{\mathrm{qual}}, r_{\mathrm{acc}}]) \cdot [r_{\mathrm{qual}}, r_{\mathrm{acc}}]
\]

最终二值化：\(r_t^{\mathrm{step}} \geq 0.5\) 时步骤标签为1（高质量），否则为0。

上述设计在程序中由数据整理组件完成。代码在构建输入序列时会把每个步骤使用 \verb|<extra_0>| 进行分割，并在这些位置贴上对应的标签。训练时，模型仅在这些位置计算损失。其他位置的标签被置为 \verb|-100|，从而被损失函数忽略。

\subsubsection{轨迹级奖励的构成}

轨迹级奖励反映整个推理过程的整体质量。该奖励不关注单个步骤，而关注”方向是否正确””知识是否覆盖完整”等更宏观的问题。轨迹级奖励由两个部分构成：最终答案的正确性与知识覆盖度。

最终答案的正确性直接从输入数据的is\_valid字段获得，无需LLM调用，定义为 \(r_{\mathrm{out}} \in \{0,1\}\)。知识覆盖度则反映推理过程是否引用了任务所需的领域知识，由LLM判断覆盖指标比例：
\[
r_{\mathrm{cover}} = \frac{\text{覆盖的指标数}}{\text{总指标数}}
\]

两者整合后的轨迹级奖励为：
\[
R_{\mathrm{traj}} = r_{\mathrm{out}} + \eta \cdot r_{\mathrm{cover}}, \quad \eta = 1.5
\]

二值化：\(R_{\mathrm{traj}} \geq 1.25\) 时轨迹标签为1（高质量），否则为0。

程序中，轨迹级标签会附着在 \verb|<extra_1>| 对应的位置。数据整理函数检查该标记是否出现在序列尾部，如果存在，则写入相应的标签。

\subsubsection{联合训练目标}

训练目标旨在让模型同时学习步骤级与轨迹级两个任务。最终的联合损失为步骤级与轨迹级交叉熵的等权求和，仅在特殊符号位置计算：

\[
L_{\text{total}} = \sum_{t \in \mathcal{P}_{\text{step}}} \text{CE}(R_\phi(h_t), y_t^{\text{step}}) + \sum_{t \in \mathcal{P}_{\text{traj}}} \text{CE}(R_\phi(h_t), y_t^{\text{traj}})
\]

其中 $\mathcal{P}_{\text{step}}$ 为所有 \texttt{<extra\_0>} 位置集合，$\mathcal{P}_{\text{traj}}$ 为 \texttt{<extra\_1>} 位置集合，其余位置的 label 置为 $-100$ 被损失函数忽略。两类损失等权合并，无需引入额外的平衡超参数。

在模型实现中，该联合损失通过 HuggingFace 提供的训练器自动完成。训练器在每一步都会从数据整理器中读取带标签的输入，并在指定位置计算损失。由于代码设置了 \verb|remove_unused_columns=False|，所有标签信息都会保留，保证奖励模型在正确位置学习。

\subsubsection{迭代机制与训练稳定性}

奖励模型遵循逐步迭代的方式进行优化。每一轮训练都会重复三项操作：数据整理、标签映射与梯度更新。此机制确保每次模型都能接触到完整的奖励信息结构。

首先，数据整理器会把思维链条映射到文本序列，同时插入步骤标记与轨迹标记。随后，程序根据序列结构分配监督位置。由于模型经过重新调整词嵌入维度，其可以正确识别新增的两个标记，并在其对应位置进行预测。

第二，训练采用梯度累积策略。小批量处理与累积梯度的结合使模型能够在资源受限的条件下仍保持稳定更新。由于序列长度较长，batch size 设置为 1 每设备，搭配梯度累积步数 4（有效 batch size=8），配合 8bit AdamW 优化器与余弦学习率调度（学习率 2e-5，warmup 50 步），在保持训练稳定性的同时提升收敛效率。

第三，模型使用余弦学习率调度以平滑训练曲线，warmup 阶段避免初期梯度震荡。整个训练过程中，奖励层级的信息不断被强化，使模型逐步形成对过程质量的稳定判断。

\subsubsection{机制的整体作用}

本节所描述的奖励设计与迭代机制使模型能够逐步形成两个能力。第一，模型能够对推理链条中的关键时刻作出判断，从而具备步骤级识别能力。第二，模型在整体维度能够判断思维流程的走向，使其具备辨别推理策略合理性的能力。

通过这样的训练逻辑，奖励模型不再只是“后验判断者”，而是一个能够真正理解推理行为的过程监督器。其学习到的能力能够在推理生成、数据筛选以及强化学习中得到应用，从而为上层系统提供稳定可靠的奖励信号。

综上，奖励模型的机制以“结构清晰、信号多元、训练可控”为设计原则。代码通过标记绑定、标签映射与联合损失的形式实现了这一设计，使整个奖励生成过程具备可复现性和可迭代性，为后续优化提供坚实基础。

\section{本章小结}

本章对研究的总体设计方案进行了系统阐述，围绕数据、模型与推理机制三个维度构建了完整的技术框架。内容从研究需求出发，明确了金融推理场景的复杂性，并提出对高质量数据、过程监督以及结构化推理机制的核心要求。在此基础上，本章给出了可操作的研究路线，并围绕数据构建、模型结构、提示模板与奖励模型设计展开分层说明。

在数据层面，本章介绍了从原始任务样本到结构化推理样本的完整流程，涵盖知识抽取、路径生成、步骤级与轨迹级奖励标注等关键环节。数据结构实现了从“答案监督”向“过程监督”转变，为奖励模型训练提供了可解释、可验证的推理证据。

在模型结构层面，本章强调了基于预训练模型的扩展策略，通过引入特殊符号机制与轻量化奖励头，使模型能够识别推理片段的结构特征并完成质量判断。模型的工作机制围绕结构化输入、稀疏标签监督及长序列处理展开，确保金融推理任务在真实环境下的可执行性。

在提示模板设计层面，本章构建了多模板协同体系。从推理生成到质量评价，从知识抽取到事实核验，模板以结构化方式组织任务信息，使数据生成过程具备连贯性与一致性，为构建高质量奖励数据奠定基础。

在奖励模型设计层面，本章提出了多维度的奖励体系，并通过统一损失函数进行整合。步骤级奖励与轨迹级奖励共同作用，使模型能够同时识别局部价值与整体走向。迭代机制保证了奖励模型的稳定训练，使其能够逐步形成对金融推理过程的理解能力。

总体来看，本章通过对研究需求、整体路线、数据构建、模型机制以及奖励体系的阐述，建立了一个面向金融推理任务的系统性框架。该框架强调可解释性、结构化表达与训练可控性，为后续章节的模型训练与实证研究提供了理论基础与技术准备。


%------------------------------------------------------

\newpage
\chapter{基于财务合成数据的过程奖励模型}
\section{数据集构建}
\subsection{财务数据合成方法}

本研究的财务数据合成基于Wind/Choice上市公司财务数据库，整体流程分为两个阶段：第一阶段为QA问答对合成，第二阶段为评分标注。

\subsubsection{源数据结构与评估维度}

从Wind/Choice提取的每条财务快照包含证券代码、公司名称、统计截止日期及financial\_data。financial\_data包含五类指标：

\begin{itemize}
    \item \textbf{核心指标}：该评估维度下最重要的2--3个指标；
    \item \textbf{补充指标}：辅助参考的4--5个扩展指标；
    \item \textbf{原始财务字段}：从财报直接提取的原始科目值；
    \item \textbf{历史指标}：近2--3年同口径数据，部分早期指标可为空；
    \item \textbf{趋势特征}：每个指标的当前值、上期值、较上期变动额/率、近3年历史均值。
\end{itemize}

数据覆盖五个评估维度：盈利能力（核心指标：销售毛利率、扣非净利润）、偿债能力（资产负债率、利息保障倍数、速动比率）、营运能力（应收账款周转天数、存货周转天数、其他应收款占比）、结构性风险（存贷双高-货币资金、存贷双高-有息负债、商誉/净资产）、现金流质量（经营现金流净额CFO、盈利现金比率、自由现金流FCF）。

\subsubsection{第一阶段：QA问答对合成}

采用Proposer→Solver→Validator三Agent协作管道，在LangGraph状态图上循环迭代合成问答对。每次迭代难度递增（0.3→1.0，步长0.1），对应从单一指标分析到综合评估的递进策略。

\textbf{Proposer（提议者）}：基于financial\_data生成分析问题、分步分析过程和参考答案。系统提示词定义角色为”信贷风险评估专家”。输出包含”问题”、”分析过程”（分步骤）和”分析结论”。

\textbf{Solver（求解者）}：独立对同一问题进行推理，按正负样本比例（正样本70\%，负样本30\%）分叉为两条路径：
\begin{itemize}
    \item 正样本：正常分析，如实回答；
    \item 负样本：注入系统性错误，包含七种错误类型——计算错误、公式引用错误、指标理解错误、逻辑跳跃、因果倒置、数据误读、趋势误判。错误要求自然、专业、有迷惑性。
\end{itemize}

\textbf{Validator（验证者）}：比较参考答案与预测答案。正样本要求两者语义等价，负样本要求存在实质性差异且错误具有迷惑性。验证通过的QA对写入历史，进入下一轮迭代。

\subsubsection{第二阶段：评分标注}

对每条推理轨迹进行步骤级与轨迹级两级评分标注，生成可监督的二值标签。

\textbf{步骤级评分}：对推理轨迹中的每一步，通过LLM并发的三个维度的评判——定性评分\(r_{\mathrm{qual}}\)（三子维度0--1均值）、过程正确性\(r_{\mathrm{procedural}}\)（二值）、事实准确性\(r_{\mathrm{factual}}\)（二值）——经Softmax动态加权聚合后以0.5阈值二值化，输出步骤级0/1标签。详细公式见第3章§3.4.3。

\textbf{轨迹级评分}：由结果正确性\(r_{\mathrm{out}}\)（直接取is\_valid字段，无需LLM）与知识覆盖度\(r_{\mathrm{cover}}\)（LLM判断覆盖指标比例）加权融合，\(\eta=1.5\)，以1.25阈值二值化。详细公式见第3章§3.4.4。

\subsubsection{输出数据格式}

评分标注完成后，每条QA对输出为JSONL格式的训练样本：

\begin{verbatim}
{
  “question”: “问题文本”,
  “knowledge_items”: {“指标名”: 数值, ...},
  “steps”: {“Step 1”: “步骤内容”, “Step 2”: “...”, ...},
  “final_answer”: “最终答案”,
  “step_labels”: [1, 0, 1, ...],
  “trajectory_label”: 1
}
\end{verbatim}

步骤分隔符\verb|<extra_0>|与轨迹结束符\verb|<extra_1>|在合成阶段即嵌入轨迹文本中，训练时模型在这些特殊token位置计算交叉熵损失，实现奖励标签的精确对齐。

\subsubsection{数据质量控制}

合成后执行自动校验，包括：结构一致性检查（确认各字段完备）、特殊符号检查（确认所有<extra\_0>和<extra\_1>存在且位置正确）、最大长度检查（超出8192 tokens则截断或重新生成）。校验通过后以UTF-8编码写入JSONL文件。

\subsection{注意力权重微调策略}

在本研究的过程奖励模型训练中，我们并未采用传统的全参数更新方式。模型主干在加载预训练权重后保持冻结，仅保留奖励头（self.score）进行全参数训练。为了在有限的计算资源下提升模型对财务轨迹的理解能力，我们在主干模型外侧引入 LoRA（Low-Rank Adaptation）结构，并将其作为唯一可更新的注意力权重调节路径。该策略避免了大规模参数更新带来的不稳定，同时强化模型在财务语境下的局部推理能力。下文对其整体设计逻辑、注意力调节思路以及完整的操作流程进行分层叙述。

\subsubsection{总体思路}

该策略的核心目标是通过低秩矩阵重构的方式，对预训练模型的注意力权重施加细粒度扰动，使其在评估财务类推理轨迹时能够形成更加明显的注意力偏置。与传统参数高维更新不同，LoRA 对注意力投影矩阵（q\_proj, k\_proj, v\_proj, o\_proj）及前馈网络投影矩阵（gate\_proj, up\_proj, down\_proj）共七个模块引入低秩分解，配置 rank=16, alpha=32，仅训练 0.75\% 参数。

在冻结主干模型的前提下，LoRA 结构承担“对齐财务推理表达”的全部责任。其核心思想可以表示为：
\[
W' = W + \Delta W,\qquad \Delta W = BA,
\]
其中 $W$ 为预训练权重，$A$ 与 $B$ 为 LoRA 中的可学习低秩矩阵，$r=\text{rank}(A)=\text{rank}(B)$ 通常远小于 $W$ 的维度。该形式保证了更新的稀疏性与可控性。

\subsubsection{注意力结构的适配逻辑}

在 Fin-PRM 的任务定义中，模型需同时处理问题文本、推理轨迹、步骤分隔符以及最终答案。因此，相较于自然语言生成任务，本任务的注意力结构呈现以下特点：

(1) 模型需识别“步骤边界符号”的位置关系，并在步骤分隔处对齐 reward head 的输出。  
(2) 模型需在多个 reasoning 片段中开展选择性关注，重点捕捉逻辑跳跃点、知识引用点以及错误推理步骤。  
(3) 模型需同时构成对 trajectory-level 与 step-level 的双尺度注意力表达。

因此，我们在注意力模块中插入 LoRA 结构，通过低秩调整引导预训练模型的注意力权重在以下区域形成优先聚焦：

\[
\mathcal{A}_{\text{focus}} = \{ \text{步骤分隔符 } (<extra_0>, <extra_1>),\ \text{概念定义区域},\ \text{推理跳转处},\ \text{数值计算片段} \}.
\]

由于 reward head 的梯度将通过 LoRA 结构向注意力模块传播，该注意力调节策略能使模型在 reward 预测中建立稳定的区域性偏置。

\subsubsection{模型包装与 LoRA 植入步骤}

整个训练流程可被理解为“加载---冻结---包装---替换---训练”的序列式机制。尽管具体代码已在前述系统实现中给出，但为了便于复现实验，本小节以操作步骤的方式全面展开。

\paragraph{加载预训练主干模型}

首先加载预训练的 Qwen 模型，如图 \ref{fig:Qwen3 Transformer Architecture}，并保持其所有 Transformer block 处于可读状态。在这一阶段不进行任何权重替换，仅确保结构正确构建。可记为：

\[
\theta_{\text{base}} = \text{Load}(M_{\text{pretrained}})
\]

其中 $\theta_{\text{base}}$ 包含全部注意力矩阵、前馈网络以及嵌入层参数。
\begin{figure}[H]
    \centering
    \includegraphics[height=0.8\textwidth]{pic/Qwen3 Transformer Architecture.png}
    \caption{Qwen-3 Transformer 架构示意图。\textbf{MHA Internal}: RoPE (32K) 位置编码。注意力计算: $Q = XW_Q + \text{RoPE}$, $K = XW_K + \text{RoPE}$, $V = XW_V$, $\text{Attn} = \text{softmax}(QK^\top/\sqrt{d})V$。后4层采用局部注意力，$QK^\top$ 计算限制在窗口内。Gated FFN: $\text{gate} = \text{SiLU}(XW_g)$, $\text{up} = XW_u$, $\text{down} = (\text{gate} \odot \text{up})W_d$。多头: 32 heads (4096 → 32 × 128)。}
    \label{fig:Qwen3 Transformer Architecture}
\end{figure}

\paragraph{冻结原始参数}

为了确保训练过程中梯度不会扩散到主干参数，我们将 $\theta_{\text{base}}$ 全部冻结：

\[
\nabla \theta_{\text{base}} = 0.
\]

这样做的目的在于确保模型不会破坏通用语义能力，仅在 LoRA 路径和奖励头上做适配性更新。

\paragraph{替换目标模块为 LoRA 包裹形式}

对七个目标模块（q\_proj, k\_proj, v\_proj, o\_proj, gate\_proj, up\_proj, down\_proj）注入低秩结构，rank=16, alpha=32，使更新矩阵具有如下形式：

\[
W^{\text{attn}}_{\text{new}} = W^{\text{attn}} + \alpha BA ,
\]

其中 $\alpha$ 为缩放因子，用于控制 LoRA 更新对原始模型的扰动幅度。
整个替换步骤可分解为三个动作：

- 定位目标层：遍历所有 Transformer block 内的注意力模块及前馈网络模块；
- 注入低秩分解矩阵：对每个目标权重插入 $(A,B)$；
- 绑定参数更新路径：确保新的可训练参数被优化器识别。

此步骤结束后，模型的可更新参数集合变为：

\[
\theta_{\text{train}} = \{ \theta_{\text{LoRA}}, \theta_{\text{reward-head}}\}.
\]

\paragraph{保留 reward head（self.score）为全参数训练}

reward head 为2层MLP（Linear→ReLU→Linear→2分类logits），将特殊符号位置的隐藏状态映射为二分类质量评分，是唯一需要全维度更新的结构。通过保留其全参数训练，可以确保 reward head 有足够的自由度表达财务轨迹的逐步奖励函数。

其输出形式保持为：

\[
\hat{r}_t = f_{\text{score}}(h_t),
\]

其中 $h_t$ 为经过 LoRA 权重调节后的隐藏状态，$\hat{r}_t$ 为2维logits，取argmax转为0/1标签。

\paragraph{接入数据构造器与分隔符标注机制}

由于数据中含有步骤分隔符 $<extra_0>$、轨迹结束符 $<extra_1>$，所有标签需沿这些标记对齐，因此训练数据在进入模型前需满足：

\[
x_{\text{encoded}} = \text{Tokenizer}( \text{trace} \oplus <extra_1> ).
\]

分隔符的 attention 权重会通过 LoRA 路径被重塑，使模型在 reward 输出处形成强显著性。

\paragraph{梯度仅从 reward head 流向 LoRA 路径}

由于原始主干模型冻结，训练过程中的梯度传播链路为：

\[
\nabla L \rightarrow \theta_{\text{reward-head}} \rightarrow \theta_{\text{LoRA}},
\]

且不会更新 $\theta_{\text{base}}$。这使训练高度稳定，并保证模型的更新集中于步骤逻辑判断能力、轨迹一致性判断能力、财务概念引用的注意力表达。

\subsubsection{注意力微调的核心机制}

在 reward 学习过程中，不同类型的监督信号通过 LoRA 路径作用于注意力模块。步骤级标签传递至评分头，引导注意力在每次出现 \texttt{<extra\_0>} 时聚焦步骤分隔位置；轨迹级标签增强模型在整条推理链关键节点的关注度；同时，包含财务知识的 token 获得更强的注意力响应，从而提升模型对概念使用准确性的判断能力。

训练过程中，梯度自适应缩放防止过度扰动主干表达，步骤密集段获得跨句的注意力链接强化，奖励头的分类损失通过 LoRA 路径引导注意力模块形成对推理轨迹质量的稳定判断。

\subsubsection{整体机制总结}

注意力微调策略通过 LoRA 注入 $\rightarrow$ 注意力扰动 $\rightarrow$ 评分头判定 $\rightarrow$ 损失反传 $\rightarrow$ LoRA 更新的闭环，将财务推理 reward 学习转化为对注意力模式的结构化重塑。该策略不改变预训练模型的底层语义能力，而是在其上叠加独立的财务轨迹评估能力，是在资源有限环境下构建领域 PRM 的关键技术手段。


\section{过程奖励模型设计}

在前述数据构建与注意力微调策略的基础上，本节正式定义过程奖励模型的完整架构与训练机制。PRM的核心任务是：给定一个问题及其推理轨迹，在每一步骤边界（\texttt{<extra\_0>}位置）输出0/1质量判断，并在轨迹终点（\texttt{<extra\_1>}位置）输出整体质量评估。模型以冻结的Qwen3-8B为主干，通过LoRA低秩适配与轻量化奖励头实现过程监督，仅训练0.75\%参数。

在同期工作中，Zhu等提出的Fin-PRM\upcite{zhuFinPRMDomainSpecializedProcess2026}同样面向金融推理构建领域专用PRM，采用步骤级与轨迹级联合建模，并在CFLUE和FinQA上验证了效果。与之相比，本工作的差异化在于：(1)数据层面，通过三Agent协作管道配合七种系统性错误注入合成训练数据，而Fin-PRM依赖MC rollout自动标注；(2)训练层面，采用LoRA参数高效微调（仅训练0.75\%参数）而非全参数训练；(3)对齐层面，基于PRM构建偏好数据并采用DPO离线对齐，而非Best-of-N推理与RL奖励塑形。

\subsection{注意力权重微调机制}

注意力权重微调是PRM实现步骤级过程监督的关键机制。其核心思想是：在冻结的主干模型之上，通过LoRA低秩矩阵构建唯一的可训练注意力调节路径，使模型在评估金融推理轨迹时形成结构化的注意力偏置，同时保持预训练语义能力不受破坏。

\subsubsection{机制整体架构}

注意力微调机制由三个组件协同构成：

\begin{enumerate}
    \item \textbf{LoRA低秩适配层}：对注意力投影矩阵（q\_proj, k\_proj, v\_proj, o\_proj）及前馈网络投影矩阵（gate\_proj, up\_proj, down\_proj）共七个模块注入低秩分解矩阵 $A \in \mathbb{R}^{d_{\text{out}} \times r}$、$B \in \mathbb{R}^{r \times d_{\text{in}}}$（rank=16, $\alpha$=32），使权重更新 $\Delta W = \alpha BA$ 仅涉及 $r(d_{\text{in}} + d_{\text{out}})$ 个参数。

    \item \textbf{奖励头（Reward Head）}：一个轻量化2层MLP（Linear $\rightarrow$ ReLU $\rightarrow$ Linear $\rightarrow$ 2分类logits），读取Transformer在特殊符号位置的隐藏状态 $h_t$，映射为二分类质量评分 $\hat{r}_t = f_{\text{score}}(h_t)$。奖励头是唯一进行全参数训练的模块。

    \item \textbf{特殊符号标记}：\texttt{<extra\_0>}标记步骤边界，\texttt{<extra\_1>}标记轨迹终点。这两个特殊token在训练时嵌入输入序列，其对应的隐藏状态分别承载步骤级和轨迹级的质量信息。
\end{enumerate}

\subsubsection{信号流动与梯度传播}

在PRM训练过程中，信号沿以下闭环流动：

\[
\text{LoRA注入} \rightarrow \text{注意力扰动} \rightarrow \text{奖励头评分} \rightarrow \text{损失反传} \rightarrow \text{LoRA更新}
\]

具体而言：（1）LoRA低秩矩阵对外部输入产生细粒度注意力扰动，使模型在步骤分隔符、金融概念术语、逻辑转折点等关键区域形成优先聚焦；（2）经LoRA调节后的隐藏状态 $h_t$ 流入奖励头，产生步骤级或轨迹级的0/1预测；（3）交叉熵损失仅计算\texttt{<extra\_0>}和\texttt{<extra\_1>}位置的预测误差；（4）梯度从奖励头反向传播，经LoRA路径更新低秩矩阵参数。由于主干模型参数冻结（$\nabla \theta_{\text{base}} = 0$），梯度传播链路为：

\[
\nabla \mathcal{L} \rightarrow \theta_{\text{reward-head}} \rightarrow \theta_{\text{LoRA}}
\]

该设计的核心优势在于：主干模型的通用语义能力完整保留，LoRA承担全部的领域适配责任，奖励头负责将注意力增强后的表示转化为质量判断。三者分工明确，训练高度稳定。

\subsubsection{注意力偏置的形成}

不同来源的监督信号通过LoRA路径差异化地塑造注意力模式：

\begin{itemize}
    \item \textbf{步骤级标签}引导注意力在每次出现\texttt{<extra\_0>}时聚焦该步骤边界，使奖励头能够基于当前步骤的上下文信息做出质量判断；
    \item \textbf{轨迹级标签}增强模型在整条推理链关键节点（如最终结论之前的推导步骤）的关注度，使\texttt{<extra\_1>}位置的表示能够聚合全轨迹信息；
    \item \textbf{财务知识token}（如指标名称、数值、会计科目）获得更强的注意力响应，提升模型对概念使用准确性的判断能力。
\end{itemize}

训练过程中，梯度自适应缩放防止过度扰动主干表达，步骤密集段获得跨句的注意力链接强化，使PRM逐步形成对推理轨迹质量的稳定判断能力。

\subsection{损失函数与优化目标}

本研究的PRM训练采用基于特殊符号位置的稀疏交叉熵损失，训练目标旨在让模型同时学习步骤级与轨迹级两个分类任务。

\subsubsection{损失函数设计}

PRM的训练损失仅计算\texttt{<extra\_0>}和\texttt{<extra\_1>}位置的交叉熵，其余位置（包括自然语言文本、数字、普通标点等）的label统一置为$-100$，被PyTorch的CrossEntropyLoss自动忽略。最终的联合损失定义为：

\[
\mathcal{L}_{\text{PRM}} = \sum_{t \in \mathcal{P}_{\text{step}}} \text{CE}(\hat{y}_t, y_t^{\text{step}}) + \sum_{t \in \mathcal{P}_{\text{traj}}} \text{CE}(\hat{y}_t, y_t^{\text{traj}})
\]

其中 $\mathcal{P}_{\text{step}}$ 为所有\texttt{<extra\_0>}位置集合（每个位置对应一个推理步骤），$\mathcal{P}_{\text{traj}}$ 为\texttt{<extra\_1>}位置集合（每条轨迹一个），$\hat{y}_t$ 为模型输出的2维logits，$y_t \in \{0, 1\}$ 为对应的二值标签。两类损失等权求和，无需引入额外的平衡超参数。

该设计的核心考量在于：(1) 稀疏监督使模型专注于步骤质量判断，而非文本生成任务；(2) 等权合并避免了$\lambda$调优的额外开销，因为步骤级和轨迹级标签在数据中的分布已自然平衡（步骤级样本远多于轨迹级，但每个位置的损失贡献均等）；(3) 仅对特殊符号位置计算损失，使梯度信号精准作用于奖励头和LoRA路径。

\subsubsection{训练优化配置}

模型训练采用以下优化策略以确保在小批量、长序列条件下的稳定收敛：

\begin{table}[h]
\centering
\caption{PRM训练优化配置}
\label{tab:prm-opt}
\begin{tabular}{lcl}
\toprule
配置项 & 取值 & 说明 \\
\midrule
优化器 & 8-bit AdamW (paged) & 降低显存占用 \\
学习率 & $2\times10^{-5}$ & 较小学习率适配LoRA微调 \\
学习率调度 & 余弦退火 & 平滑训练曲线 \\
预热步数 & 50 & 避免初期梯度震荡 \\
训练精度 & bf16 & 平衡精度与效率 \\
每设备批次 & 1 & 适配长序列（3072 tokens） \\
梯度累积步数 & 4 & 等效批次大小 = 8 \\
最大序列长度 & 3072 tokens & 覆盖完整推理轨迹 \\
\bottomrule
\end{tabular}
\end{table}

训练过程中采用梯度累积策略：由于单条金融推理样本的序列长度较长（含问题、知识项、三步推理及标注），单设备仅能容纳batch\_size=1，通过累积4步梯度后统一更新，实现等效batch\_size=8的训练效果。余弦学习率调度配合50步预热，在训练初期平缓提升学习率以避免随机初始化奖励头带来的梯度震荡，随后余弦衰减至接近零，确保后期精细收敛。

训练共进行5个epoch，约1,250步。实验表明（详见§5.4.1），模型在Epoch 4--5已达到收敛饱和：loss标准差从Epoch 1的5.1降至Epoch 5的0.36，最终loss为0.153（较初始16.73下降99.1\%），5 epoch是充分的训练设置。

\section{实验验证}
\subsection{实验环境设置}

PRM训练与DPO训练均在同一台服务器上完成，核心软硬件配置如表~\ref{tab:exp-env}所示。

\begin{table}[h]
\centering
\caption{实验环境配置}
\label{tab:exp-env}
\begin{tabular}{lcc}
\toprule
\multirow{2}{*}{配置项} & \multicolumn{2}{c}{训练阶段} \\
\cmidrule(lr){2-3}
& PRM 过程奖励模型 & DPO 策略模型 \\
\midrule
GPU & \multicolumn{2}{c}{NVIDIA RTX PRO 6000 Blackwell, 98GB} \\
基座模型 & \multicolumn{2}{c}{Qwen3-8B} \\
微调方式 & LoRA (rank=16, $\alpha$=32) & LoRA (rank=16, $\alpha$=32) \\
训练精度 & bf16 & bf16 \\
优化器 & paged\_adamw\_8bit & paged\_adamw\_8bit \\
学习率 & $2\times10^{-5}$ & $5\times10^{-6}$ \\
批次大小 & 8 (累积梯度) & 16 \\
训练轮数 & 5 epochs & 1 epoch \\
训练时长 & $\sim$2.5h (9,143s) & $\sim$8.9min (531s) \\
框架 & HuggingFace Trainer + PEFT & TRL DPOTrainer + PEFT \\
\bottomrule
\end{tabular}
\end{table}

PRM训练与DPO训练采用差异化学习率策略：PRM作为评分模型需要充分学习步骤质量判断，采用$2\times10^{-5}$；DPO训练中LoRA仅需微调策略模型的偏好方向，采用更保守的$5\times10^{-6}$以保持基座模型的推理能力。

\subsection{实验数据集描述}

本实验使用的数据集涵盖PRM训练与DPO偏好训练两个阶段，各数据集的样本数量与用途如表~\ref{tab:dataset}所示。

\begin{table}[h]
\centering
\caption{实验数据集概况}
\label{tab:dataset}
\begin{tabular}{lcc}
\toprule
数据集 & 样本数 & 用途 \\
\midrule
PRM Train & 1,995 & PRM 过程奖励模型训练 \\
PRM Validate & 400 & PRM 训练中验证 \\
PRM Test & 200 & PRM 最终测试评估 \\
DPO Train & 2800 对 & DPO 偏好对齐训练 \\
DPO Validate & 800 对 & DPO 训练中验证 \\
DPO Test & 400 对 & DPO 最终测试评估 \\
\bottomrule
\end{tabular}
\end{table}

每条PRM样本包含一个金融分析问题、约70项财务指标（knowledge\_items）、3步推理轨迹及对应的步骤级/轨迹级二值标签。类别分布：正确步骤约占82\%，错误步骤约占18\%；正轨迹（所有步骤正确且答案正确）约占37\%，负轨迹约占63\%。DPO偏好对由PRM对同一问题的4条候选回答进行评分排序构建：PRM对每条候选输出步骤级与轨迹级评分，通过"步骤平均+加权融合"计算综合得分，chosen为最高分回答，rejected为最低分回答。

\subsection{实验评估指标}

\textbf{PRM评估指标}：本研究从分类性能与置信度校准两个维度评估PRM：

\begin{itemize}
    \item \textbf{步骤准确率}（Step Accuracy）：步骤级二分类正确率，为核心指标，目标$>$90\%；
    \item \textbf{轨迹准确率}（Trajectory Accuracy）：整条推理轨迹的标签预测正确率；
    \item \textbf{正确步骤召回率}（Correct Step Recall）：正确步骤中被PRM识别为正确的比例；
    \item \textbf{错误步骤召回率}（Error Step Recall）：错误步骤中被PRM识别为错误的比例；
    \item \textbf{置信度校准}：正确步骤与错误步骤的平均prob\_1差距，差距越大校准越好。
\end{itemize}

\textbf{DPO评估指标}：DPO策略模型的评估包括训练过程指标与偏好对齐指标，目标值参考实验方案设定：

\begin{table}[h]
\centering
\caption{DPO策略模型评估指标与目标}
\label{tab:dpo-metrics}
\begin{tabular}{lcc}
\toprule
维度 & 指标 & 目标 \\
\midrule
偏好对齐 & Preference Accuracy（chosen奖励$>$rejected的比例） & $>$70\% \\
推理准确率 & 最终答案正确率 & 提升10--15\% \\
策略稳定性 & DPO Loss & 稳定下降 \\
参考偏离度 & $\text{KL}(\pi_\theta \| \pi_{\text{ref}})$ & $<$5.0 \\
\bottomrule
\end{tabular}
\end{table}

\subsection{实验结果与分析}

本节以七个实验子节系统报告PRM训练、评估与DPO偏好数据构建、策略训练的完整实验结果。所有实验均基于§5.1所述环境配置，所用图表均来源于训练过程中记录的实际数据。

\subsubsection{实验一：PRM训练收敛性分析}

本实验分析PRM在5个epoch（1,250步）训练过程中的loss收敛行为，验证训练设置的充分性。

\begin{figure}[H]
    \centering
    \includegraphics[width=0.95\textwidth]{pic/01_loss_curve.png}
    \caption{PRM训练完整loss曲线（对数坐标）}
    \label{fig:loss-curve}
\end{figure}

图~\ref{fig:loss-curve}展示了完整的训练loss轨迹。初始loss高达16.73（step 10），模型处于随机猜测阶段。Epoch 1（step 0--250）loss快速下降但剧烈震荡，前期最高spike达32.95，反映模型在探索不同难度样本的步骤质量判断模式。Epoch 2--3（step 250--750）loss降至1.0附近，波动幅度逐步收敛。Epoch 4--5（step 750--1250）loss进一步降至0.02--1.6范围，趋于稳定。最终loss为0.153（step 1250），较初始下降99.1\%。

\begin{figure}[H]
    \centering
    \includegraphics[width=0.95\textwidth]{pic/02_loss_per_epoch.png}
    \caption{逐epoch loss分图}
    \label{fig:loss-per-epoch}
\end{figure}

图~\ref{fig:loss-per-epoch}将5个epoch的loss分别绘制于5个子图中，虚线标注各epoch均值。从per-epoch平均loss的演变清晰可见：Epoch 1均值5.37（极不稳定）→ Epoch 2均值3.43 → Epoch 3均值1.84 → Epoch 4均值1.52 → Epoch 5均值0.38。Epoch 2中出现一次极端spike（loss=32.95），可能由某批次样本质量异常或标注噪声引起。Epoch 5中绝大部分loss已降至1.0以下，最低至0.018。

\begin{figure}[H]
    \centering
    \includegraphics[width=0.95\textwidth]{pic/03_loss_distribution.png}
    \caption{Loss分布演变直方图}
    \label{fig:loss-dist}
\end{figure}

图~\ref{fig:loss-dist}以直方图展示loss分布从"分散+长尾"到"集中+近零"的演变过程。Epoch 1--2的loss呈右偏分布，大部分集中在0--5区间但有长尾延伸至30+；Epoch 3起分布显著左移；Epoch 5的loss高度集中在0--0.5范围。整体分布均值为2.46，分布形态的演化直观验证了训练的有效性。

\begin{figure}[H]
    \centering
    \includegraphics[width=0.95\textwidth]{pic/05_convergence_analysis.png}
    \caption{收敛性分析三合一图}
    \label{fig:convergence}
\end{figure}

图~\ref{fig:convergence}从三个维度分析收敛：（左）各epoch结束时的loss终点值——epoch 1终点1.65 → epoch 5终点0.15，下降91\%；（中）累计loss下降率——epoch 2已达36\%，epoch 3达62\%，epoch 5达93\%，表明大部分下降发生在前3个epoch；（右）各epoch的loss标准差——从epoch 1的std=5.1（极不稳定）降至epoch 5的std=0.36（高度稳定）。

\textbf{小结}：PRM在5个epoch内实现了从随机猜测到稳定判断的跨越，loss下降99.1\%。收敛分析显示，Epoch 5的loss下降和稳定性均已饱和，增加epoch带来的边际收益递减——5 epoch是充分的训练设置。

\subsubsection{实验二：PRM步骤级与轨迹级评估}

本实验在验证集（400样本）和测试集（200样本）上评估PRM的分类性能，结果如表~\ref{tab:prm-eval}所示。

\begin{table}[h]
\centering
\caption{PRM步骤级与轨迹级评估结果}
\label{tab:prm-eval}
\begin{tabular}{lccc}
\toprule
指标 & Validate & Test & 是否达标 \\
\midrule
步骤准确率 & \textbf{95.25\%} & \textbf{94.00\%} & $\checkmark$（目标$>$90\%） \\
轨迹准确率 & 72.00\% & 78.00\% & — \\
正确步骤召回率 & 97.7\% & 97.3\% & — \\
错误步骤召回率 & 84.4\% & 79.5\% & — \\
\bottomrule
\end{tabular}
\end{table}

步骤准确率在验证集和测试集上分别达到95.25\%和94.00\%，均超过90\%目标，且两集差异仅1.25个百分点，表明模型泛化良好。轨迹准确率（72\%--78\%）低于步骤级准确率，这是合理的：轨迹判断需聚合所有步骤信号，任意一步误判即导致轨迹级错误。

正确步骤召回率（97\%）显著高于错误步骤召回率（80\%），与类别不平衡（正确步骤82\% vs 错误步骤18\%）一致。这一现象可通过分析训练集中正负样本比例来解释：模型见过更多正确步骤，对其特征学习更充分。错误步骤召回率从旧版的0.9\%提升至79.5\%（详见实验四），是PRM训练策略优化的核心成果。

\begin{figure}[H]
    \centering
    \includegraphics[width=0.95\textwidth]{pic/08_loss_accuracy.png}
    \caption{Loss与步骤准确率对照}
    \label{fig:loss-acc}
\end{figure}

图~\ref{fig:loss-acc}展示了loss下降与准确率提升的正相关关系，同时也揭示了loss的边际递减效应：Epoch 1时loss约5、准确率约90\%；Epoch 4--5时loss从1.0降至0.15，准确率仅从93\%微幅提升至94\%。这表明loss并非准确率的线性代理指标，在训练后期应更关注评估指标本身而非loss的持续下降。

\subsubsection{实验三：置信度校准分析}

置信度校准是PRM可用性的关键：如果模型对所有步骤输出接近的置信度，则无法有效区分正确与错误步骤。本实验分析新版PRM（ckpt-1250, 5ep）在测试集上的置信度分布。

\begin{table}[h]
\centering
\caption{PRM置信度校准结果}
\label{tab:confidence}
\begin{tabular}{lcc}
\toprule
指标 & Validate & Test \\
\midrule
正确步骤平均prob\_1 & 0.9747 & 0.9741 \\
错误步骤平均prob\_1 & 0.1771 & 0.2408 \\
置信度差距 & \textbf{0.80} & \textbf{0.73} \\
\bottomrule
\end{tabular}
\end{table}

正确步骤的置信度高度集中在0.9--1.0区间，错误步骤集中在0.0--0.3区间，两者均值差距在验证集和测试集上分别达到0.80和0.73。这表明PRM不是简单地对所有样本输出高置信度，而是具有良好校准的判断能力。

\begin{figure}[H]
    \centering
    \includegraphics[width=0.9\textwidth]{pic/07_confidence_analysis.png}
    \caption{置信度分布与混淆矩阵}
    \label{fig:confidence}
\end{figure}

图~\ref{fig:confidence}（左）以直方图展示正确/错误步骤的prob\_1分布，黑色虚线标注0.5决策边界。正确步骤的prob\_1高度集中在0.9--1.0区间（488个中475个$\geq$0.5），错误步骤集中在0.0--0.3区间（112个中89个$<$0.5）。图~\ref{fig:confidence}（右）的混淆矩阵显示：TP=475, FN=13, FP=23, TN=89。FP=23意味着6.2\%的错误步骤被误判为正确，这些样本可能包含表面合理但存在隐蔽错误的步骤；FN=13表示2.7\%的正确步骤被误判为错误，可能涉及表述不清晰但逻辑正确的步骤。

新版PRM的置信度校准良好，正确和错误步骤的置信度分别为0.974和0.241，差距达0.73。旧版因7 epoch训练导致logits极端化，两类步骤置信度均接近1.0（0.999和0.981），几乎无区分度（详见实验四）。缩减训练epoch、避免logits极端化，是恢复置信度区分能力的关键。

\subsubsection{实验四：过拟合对比分析}

本实验是本研究最重要的实验发现。在PRM训练初期，使用7 epoch的配置获得了看似合理的步骤准确率（81.3\%），但进一步分析发现了严重的过拟合问题。

\begin{table}[h]
\centering
\caption{旧版与新版的PRM评估对比}
\label{tab:overfit}
\begin{tabular}{lccc}
\toprule
指标 & 旧版（7ep） & 新版（5ep） & 改进幅度 \\
\midrule
步骤准确率 & 81.3\% & 94.0\% & +12.7\% \\
轨迹准确率 & 46.0\% & 78.0\% & +32.0\% \\
错误步骤召回率 & \textbf{0.9\%} & \textbf{79.5\%} & \textbf{+78.6\%} \\
正确步骤置信度 & 0.999 & 0.974 & 更校准 \\
错误步骤置信度 & 0.981 & 0.241 & 更有区分度 \\
\bottomrule
\end{tabular}
\end{table}

旧版PRM的错误步骤召回率仅为0.9\%——模型几乎将所有步骤（无论正确与否）都判为正确。同时，旧版对正确和错误步骤均输出接近1.0的置信度（0.999和0.981），完全不具备区分能力。这是典型的小样本过拟合表现：模型学会的是"几乎所有步骤都是正确的"这一退化解，而非真正的步骤质量判断能力。

\begin{figure}[H]
    \centering
    \includegraphics[width=0.95\textwidth]{pic/06_eval_comparison.png}
    \caption{旧版与新版的PRM评估综合对比}
    \label{fig:eval-comp}
\end{figure}

图~\ref{fig:eval-comp}以双面板综合对比了新旧版本：（左）四项指标的柱状对比，橙色虚线标注90\%目标线——新版在步骤准确率和错误步骤召回率上均大幅领先；（右）置信度对比——旧版对两类步骤均输出$\sim$1.0的高置信度（无区分为），新版实现了正确步骤0.97、错误步骤0.24的清晰分离。

\textbf{过拟合的根因与修复}：旧版使用7个epoch训练，在小样本（1,995条）条件下模型过度记忆了训练数据的表面模式——大部分训练样本的正确步骤占比高（82\%），模型学会通过"一律判为正确"即可在训练集上获得较低loss。修复措施为：将训练epoch从7缩减至5，在loss收敛后及时停止，避免logits极端化，从而增大正负类间的区分度。

\textbf{核心贡献}：在小样本PRM训练中，epoch数是防止过拟合的关键控制变量。仅通过缩减训练epoch这一调整，错误步骤召回率即从0.9\%恢复至79.5\%，步骤准确率提升12.7个百分点。这一发现对后续同类小样本PRM训练工作具有重要参考价值。

\subsubsection{实验五：DPO偏好数据集质量分析}

本实验分析基于新版PRM构建的DPO偏好数据集的质量特征。偏好对构建方法：对每个问题生成多条候选回答，由PRM对每条候选输出步骤级（\texttt{<extra\_0>}位置）与轨迹级（\texttt{<extra\_1>}位置）评分，通过"步骤平均+加权融合"（$\text{score}=\alpha\times\text{mean}(\text{步骤级})+\beta\times\text{轨迹级}$）计算综合得分并排序，chosen为最高分回答，rejected为最低分回答。实验首先在399对规模（每题5条候选）上验证PRM修复过拟合对偏好数据质量的提升，随后将规模扩展至4000对（每题4条候选）以支撑DPO训练。

\begin{figure}[H]
    \centering
    \includegraphics[width=0.95\textwidth]{pic/14_baseline_comparison.png}
    \caption{旧版与新版PRM构建的DPO偏好数据质量对比}
    \label{fig:baseline-cmp}
\end{figure}

图~\ref{fig:baseline-cmp}给出了关键的过渡性对比。旧版PRM产出的偏好数据几乎无区分度：平均margin仅0.004，中位数margin为0.000，margin$\geq$0.2的样本仅占0.5\%。新版PRM修复过拟合后，数据质量发生根本性改善：平均margin提升87倍至0.348，中位数margin提升至0.311，margin$\geq$0.2的占比提升129倍至64.5\%。这定量验证了"先解决PRM过拟合，再构建偏好数据"策略的正确性。

\begin{figure}[H]
    \centering
    \includegraphics[width=0.95\textwidth]{pic/11_margin_distribution.png}
    \caption{DPO偏好数据三分集margin分布}
    \label{fig:margin-dist}
\end{figure}

图~\ref{fig:margin-dist}展示了Train（279对）、Validate（80对）、Test（40对）三个数据集的margin分布。分布呈类正态形态，Train集峰值在[0.30, 0.40)区间（47个，16.8\%），中位数0.31。Validate集分布更均匀，中位数0.40。三个数据集的分布形态一致，未见明显的分布偏移，说明数据划分具有良好的统计同质性。

\begin{figure}[H]
    \centering
    \includegraphics[width=0.95\textwidth]{pic/12_margin_stats.png}
    \caption{Margin统计：箱线图与阈值达标率}
    \label{fig:margin-stats}
\end{figure}

图~\ref{fig:margin-stats}（左）以箱线图展示三数据集的margin五数概括。Validate集整体最高（均值0.40），四分位距（P25--P75: 0.16--0.63）最宽，偏好信号最强。Train和Test集水平接近（均值0.33）。图~\ref{fig:margin-stats}（右）展示了不同阈值下的达标率：92\%--93\%样本margin$\geq$0.05（具有正向区分度），65\%--69\%样本margin$\geq$0.20（明确区分），20\%--39\%样本margin$\geq$0.50（高质量区分）。无margin为0.00的异常集中，说明偏好对不同质。

\begin{figure}[H]
    \centering
    \includegraphics[width=0.95\textwidth]{pic/13_score_analysis.png}
    \caption{Chosen vs Rejected分数散点图与margin排序曲线}
    \label{fig:score-analysis}
\end{figure}

图~\ref{fig:score-analysis}进一步分析了分数的分布特征。（左）散点图显示大量点集中在右上角（chosen/rejected均高分），但rejected分数广泛分布在0.0--1.0区间，形成明显的对角线偏离——高margin样本（绿色）远离对角线，低margin样本（红色）靠近对角线。（右）排序曲线显示margin从0.01平滑上升至0.96，拐点约在第100对（约35\%的样本margin$<$0.2），64.5\%的样本margin$\geq$0.2。

\begin{figure}[H]
    \centering
    \includegraphics[width=0.7\textwidth]{pic/15_overall_summary.png}
    \caption{DPO偏好数据整体质量总览}
    \label{fig:dpo-summary}
\end{figure}

图~\ref{fig:dpo-summary}以饼图总结：399对偏好数据中，26.3\%（105对）margin$\geq$0.5（高质量），73.7\% margin$<$0.5。92.2\%的样本margin$\geq$0.05，近乎所有偏好对都具有正向区分度。高质量样本约占四分之一，足以支撑有效的偏好学习。

为进一步验证偏好数据的规模化构建能力并支撑后续DPO训练，本研究将候选回答从5条调整为4条、将数据规模从399对扩充至4000对（Train 2800 + Validate 800 + Test 400），得到平均margin 0.686、中位数0.762的偏好数据集，95.4\%的样本margin$\geq$0.2。图~\ref{fig:margin-dist-4000}展示了三分集的margin分布：分布呈右偏形态，峰值位于[0.80, 0.90)区间（22.5\%），约43\%的样本margin$>$0.8，仅4.6\%的样本margin$<$0.2。图~\ref{fig:baseline-cmp-4000}对比了新旧构建的关键质量指标：平均margin从0.348提升至0.686（约2倍），中位数margin从0.311提升至0.762（约2.5倍），margin$\geq$0.2的占比从64.5\%提升至95.4\%（+31个百分点），margin$\geq$0.5的占比从23.7\%提升至77.0\%（+53个百分点）。chosen与rejected分数进一步分离（0.858 vs 0.171），偏好信号显著增强。三集统计高度一致，无拆分偏差。

\begin{figure}[H]
    \centering
    \includegraphics[width=0.95\textwidth]{pic/18_margin_distribution_4000.png}
    \caption{4000对偏好数据集三分集margin分布}
    \label{fig:margin-dist-4000}
\end{figure}

\begin{figure}[H]
    \centering
    \includegraphics[width=0.95\textwidth]{pic/21_baseline_comparison_4000.png}
    \caption{新旧DPO偏好数据构建质量对比（399对 vs 4000对）}
    \label{fig:baseline-cmp-4000}
\end{figure}

\textbf{小结}：新版PRM使DPO偏好数据的平均区分度从0.004提升至0.348（87倍）；进一步将数据规模扩展至4000对后，平均margin提升至0.686，95.4\%的样本margin$\geq$0.2，为DPO策略训练提供了高质量的偏好信号基础。

\subsubsection{实验六：DPO策略模型训练}

本实验分两步报告DPO策略模型的训练结果：首先在279对偏好数据上进行初步训练以诊断数据规模瓶颈，随后将数据规模扩充至2800对完成完整训练。

\paragraph{初步训练（279对）：数据规模瓶颈验证}

首先在279对偏好数据上进行1 epoch的初步DPO训练，结果如表~\ref{tab:dpo-results}所示。

\begin{table}[h]
\centering
\caption{DPO初步训练结果与实验方案目标对比}
\label{tab:dpo-results}
\begin{tabular}{lccc}
\toprule
指标 & 实验方案目标 & 实际结果 & 状态 \\
\midrule
训练数据量 & 5,000--10,000 对 & 279 对 & 差18--36倍 \\
Preference Accuracy & $>$70\% & 42.5\% & 低于50\%随机基线 \\
Rewards Margin & 正值 & $-$0.005至$-$0.029 & 无法区分 \\
DPO Loss & 稳定下降 & 0.693$\rightarrow$0.701 & 接近ln2$\approx$0.693 \\
训练时长 & — & 531s ($\sim$8.9min) & — \\
LoRA适配器大小 & — & 167 MB & — \\
\bottomrule
\end{tabular}
\end{table}

\begin{figure}[H]
    \centering
    \includegraphics[width=0.95\textwidth]{pic/16_dpo_training_metrics.png}
    \caption{DPO初步训练指标演进}
    \label{fig:dpo-training}
\end{figure}

图~\ref{fig:dpo-training}以四合一图展示了初步训练过程的完整指标变化：（左上）DPO loss始终在0.693--0.702窄幅波动，接近随机基线ln(2)=0.693，1个epoch下模型尚未显著偏离参考模型；（右上）Rewards Accuracy从0\%逐步提升至47.5\%，呈上升趋势但未突破50\%随机基线；（左下）Chosen与Rejected的隐式奖励均在0附近波动且差距极小；（右下）Rewards Margin全为负值（$-$0.005至$-$0.029），模型在大多数batch中仍倾向rejected回答。

\begin{figure}[H]
    \centering
    \includegraphics[width=0.9\textwidth]{pic/17_dpo_summary.png}
    \caption{DPO数据质量与训练总览}
    \label{fig:dpo-overview}
\end{figure}

图~\ref{fig:dpo-overview}（左）展示了三数据集的平均margin与margin$\geq$0.2占比对比——Validate集在两指标上均最优。（右）列出了训练关键参数与校验结果——所有训练配置均与实验方案一致（全部$\checkmark$校验通过），说明当前的瓶颈明确指向数据规模而非训练设计。

\textbf{数据规模瓶颈的定量验证}：初步训练的Loss接近随机基线、margin为负、Rewards Accuracy低于50\%——这三个现象一致指向同一结论：偏好信号数量级不足以驱动显著的偏好对齐。这不是方案设计的失败，而是数据规模的定量验证。

\paragraph{规模扩展后的完整训练（2800对）：首次成功收敛}

基于上述诊断，将偏好数据规模扩充至4000对（Train 2800 + Validate 800 + Test 400，平均margin 0.686），重新进行1 epoch DPO训练。训练配置与实验方案完全一致（LoRA rank=16、$\alpha$=32、学习率5e-6、$\beta$=0.1、有效batch 16、最大序列长度3072），关键训练指标的逐步演进如表~\ref{tab:dpo-final-steps}所示。

\begin{table}[h]
\centering
\caption{DPO完整训练指标演进（2800对，1 epoch）}
\label{tab:dpo-final-steps}
\begin{tabular}{lcccccc}
\toprule
Step & Loss & Rewards Acc & Rewards Margin & Chosen Reward & Rejected Reward \\
\midrule
1 & 0.693 & 0.000 & 0.000 & 0.000 & 0.000 \\
30 & 0.663 & 0.637 & 0.070 & 0.155 & 0.085 \\
60 & 0.284 & 0.994 & 1.283 & 2.204 & 0.922 \\
90 & 0.128 & 0.944 & 3.899 & 6.297 & 2.398 \\
130 & 0.109 & 0.975 & 5.122 & 7.585 & 2.463 \\
170 & 0.084 & 0.969 & 5.250 & 7.736 & 2.486 \\
\bottomrule
\end{tabular}
\end{table}

在留出的验证集与测试集上评估DPO隐式奖励的偏好准确率（即chosen隐式奖励高于rejected隐式奖励的比例），结果如表~\ref{tab:dpo-final-eval}所示。

\begin{table}[h]
\centering
\caption{DPO策略模型留出集评估结果}
\label{tab:dpo-final-eval}
\begin{tabular}{lcccc}
\toprule
数据集 & 样本数 & 偏好准确率 & 平均margin & chosen奖励 \\
\midrule
Validate & 800 & \textbf{99.25\%}（794/800） & 5.92 & 7.96 \\
Test & 400 & \textbf{99.50\%}（398/400） & 6.20 & 8.10 \\
\bottomrule
\end{tabular}
\end{table}

训练loss从随机基线ln2$\approx$0.693单调下降至末步0.084（全epoch平均0.275），Rewards Accuracy在前60步即突破90\%并稳定于95\%--99\%，Rewards Margin从0稳定爬升至约5.0，chosen/rejected隐式奖励显著分离（约7.7 vs 2.5）。在留出的验证集与测试集上，偏好准确率分别达99.25\%与99.50\%，远超实验方案70\%的目标，且两集表现一致、无过拟合迹象。表~\ref{tab:dpo-compare}对比了新旧两次训练的关键指标。

\begin{table}[h]
\centering
\caption{DPO新旧训练结果对比}
\label{tab:dpo-compare}
\begin{tabular}{lcc}
\toprule
指标 & 旧训练（279对） & 新训练（2800对） \\
\midrule
训练样本 & 279 & 2800（10$\times$） \\
全epoch平均Loss & 0.701 & \textbf{0.275} \\
Rewards Accuracy & 47.5\% & \textbf{$\sim$97\%} \\
Rewards Margin & 负（$-$0.03） & \textbf{$\sim$5.0} \\
偏好准确率（held-out） & 42.5\% & \textbf{99.25\% / 99.50\%} \\
\bottomrule
\end{tabular}
\end{table}

\textbf{小结}：将训练数据从279对扩充至2800对（10倍）后，DPO训练首次成功收敛，偏好准确率从42.5\%提升至99.25\%/99.50\%，定量验证了"数据规模是偏好对齐关键瓶颈"的判断，证明PRM驱动的DPO离线偏好对齐路线在金融推理场景下切实可行。

\subsubsection{实验七：偏好构造方式与超参数消融}

为系统评估偏好构造方式与超参数对DPO偏好对齐的影响，本研究在4000对偏好数据基础上开展了三组消融实验：偏好构造方式消融、$\beta$参数消融与训练epoch消融。

\textbf{偏好构造方式消融}：对比四种排序方式——仅轨迹级（trajectory）、仅步骤平均（step\_mean）、静态权重$\alpha=0.5$（fusion\_static）与加权融合$\alpha=0.7,\beta=0.3$（fusion\_weighted）。各方式在同一份PRM打分缓存上离线计算综合得分并构建偏好对，结果如表~\ref{tab:abl-ranking}所示。

\begin{table}[h]
\centering
\caption{偏好构造方式消融结果（$\beta$=0.1，1 epoch）}
\label{tab:abl-ranking}
\begin{tabular}{lccc}
\toprule
排序方式 & 验证集偏好准确率 & 测试集偏好准确率 & 平均margin \\
\midrule
仅轨迹级（基线） & 96.88\% & 96.50\% & 4.94 \\
仅步骤平均 & \textbf{99.25\%} & \textbf{99.50\%} & \textbf{5.92} \\
静态权重$\alpha$=0.5 & 98.88\% & 98.50\% & 5.72 \\
加权融合$\alpha$=0.7/$\beta$=0.3 & 98.88\% & 99.00\% & 5.73 \\
\bottomrule
\end{tabular}
\end{table}

结果表明：仅步骤平均排序取得最优偏好准确率（99.25\%/99.50\%）与最大margin（5.92），显著优于仅轨迹级基线（96.88\%/96.50\%）。加权融合（98.88\%/99.00\%）未带来额外增益，反而略低于纯步骤平均，说明步骤级评分（PRM步骤准确率94\%）比轨迹级评分（78\%）更可靠，偏好信号的核心来源是步骤级评分本身，而非融合方式。

\textbf{$\beta$参数消融}：在加权融合排序上对比$\beta$=0.01/0.05/0.10/0.50，结果如表~\ref{tab:abl-beta}所示。偏好准确率对$\beta$不敏感（稳定在98.88\%/99.50\%），而margin随$\beta$单调放大（0.94$\rightarrow$11.21），说明$\beta$控制的是偏离参考模型的程度而非偏好对齐准确率，$\beta$=0.10仍是margin与稳定性的合理折中。

\begin{table}[h]
\centering
\caption{$\beta$参数消融结果（加权融合排序，1 epoch）}
\label{tab:abl-beta}
\begin{tabular}{lccc}
\toprule
$\beta$ & 验证集偏好准确率 & 测试集偏好准确率 & 平均margin \\
\midrule
0.01 & 98.88\% & 99.50\% & 0.94 \\
0.05 & 98.88\% & 99.50\% & 3.78 \\
0.10 & 98.88\% & 99.00\% & 5.73 \\
0.50 & 98.88\% & 99.50\% & 11.21 \\
\bottomrule
\end{tabular}
\end{table}

\textbf{epoch消融}：在加权融合排序上对比1/2/3 epoch，结果如表~\ref{tab:abl-epoch}所示。偏好准确率随epoch单调微升（98.88\%$\rightarrow$99.25\%），margin持续放大（5.73$\rightarrow$8.44），未观察到reward hacking现象。

\begin{table}[h]
\centering
\caption{epoch消融结果（加权融合排序，$\beta$=0.1）}
\label{tab:abl-epoch}
\begin{tabular}{lccc}
\toprule
epoch & 验证集偏好准确率 & 测试集偏好准确率 & 平均margin \\
\midrule
1 & 98.88\% & 99.00\% & 5.73 \\
2 & 99.25\% & 99.50\% & 7.71 \\
3 & 99.25\% & 99.50\% & 8.44 \\
\bottomrule
\end{tabular}
\end{table}

\textbf{小结}：偏好构造方式消融表明仅步骤平均排序最优（99.25\%/99.50\%），验证了步骤级评分是偏好信号的核心来源；$\beta$与epoch消融均复现了"$\beta$只影响margin、epoch单调微升无过拟合"的结论。据此，最终DPO策略模型采用仅步骤平均排序方式构建偏好数据。

\subsection{消融实验}

消融实验旨在系统评估各关键组件对PRM与DPO性能的贡献，涵盖PRM侧（打分温度、训练epoch、训练数据规模）与DPO侧（偏好构造方式、$\beta$参数、训练epoch）两组实验。DPO侧消融结果已在上文实验七报告，本小节重点报告PRM侧三组消融结果。

\textbf{打分温度消融}：在checkpoint-1250（5 epoch）上对比打分温度T=1.0/1.5/2.0/2.5，结果如表~\ref{tab:abl-temp}所示。步骤准确率与正确/错误步骤召回率在T=1.0--2.5完全一致（94.00\%/97.34\%/79.46\%），说明以prob\_1$\geq$0.5为阈值的分类决策具有T-不变性；但置信度差距随T增大单调缩小（0.73$\rightarrow$0.63），T=1.0的置信度区分度最优。

\begin{table}[h]
\centering
\caption{PRM打分温度消融结果（checkpoint-1250，测试集）}
\label{tab:abl-temp}
\begin{tabular}{lcccc}
\toprule
温度T & 步骤准确率 & 正确步骤召回率 & 错误步骤召回率 & 置信度差距 \\
\midrule
1.0 & 94.00\% & 97.34\% & 79.46\% & 0.733 \\
1.5 & 94.00\% & 97.34\% & 79.46\% & 0.704 \\
2.0 & 94.00\% & 97.34\% & 79.46\% & 0.669 \\
2.5 & 94.00\% & 97.34\% & 79.46\% & 0.632 \\
\bottomrule
\end{tabular}
\end{table}

\textbf{训练epoch消融}：固定T=1.0与1995样本，对比epoch=3/5/7，结果如表~\ref{tab:abl-epoch-prm}所示。步骤准确率与错误召回率随epoch先升后崩——3 epoch欠拟合（92.00\%/63.4\%）、5 epoch最优（94.00\%/79.5\%）、7 epoch过拟合崩坏（81.3\%/0.9\%），模型退化为几乎将所有步骤判为正确。

\begin{table}[h]
\centering
\caption{PRM训练epoch消融结果（测试集）}
\label{tab:abl-epoch-prm}
\begin{tabular}{lccc}
\toprule
epoch & 步骤准确率 & 错误步骤召回率 & 置信度差距 \\
\midrule
3 & 92.00\% & 63.39\% & 0.594 \\
5 & \textbf{94.00\%} & \textbf{79.46\%} & \textbf{0.733} \\
7 & 81.30\% & 0.90\% & 0.02 \\
\bottomrule
\end{tabular}
\end{table}

\textbf{训练数据规模消融}：固定T=1.0与5 epoch，对比训练样本500/1000/1995，结果如表~\ref{tab:abl-data}所示。步骤准确率与错误召回率随数据量单调提升（86.3\%$\rightarrow$94.0\%，33.0\%$\rightarrow$79.5\%），500样本时呈现典型"小样本过拟合"退化（正确召回仍98.6\%但错误召回仅33.0\%），定量印证了错误步骤识别能力对数据量高度敏感。

\begin{table}[h]
\centering
\caption{PRM训练数据规模消融结果（测试集）}
\label{tab:abl-data}
\begin{tabular}{lccc}
\toprule
样本数 & 步骤准确率 & 错误步骤召回率 & 置信度差距 \\
\midrule
500 & 86.33\% & 33.04\% & 0.323 \\
1000 & 91.83\% & 63.39\% & 0.580 \\
1995 & \textbf{94.00\%} & \textbf{79.46\%} & \textbf{0.733} \\
\bottomrule
\end{tabular}
\end{table}

\textbf{小结}：PRM侧三组消融表明——打分温度对分类结果无影响（仅影响置信度）、epoch=5为最优（3欠拟合、7过拟合）、错误步骤识别能力随数据量单调提升。结合实验七的DPO侧消融，偏好构造方式（仅步骤平均）、$\beta$（0.10）与epoch（1）等关键配置均已通过消融实验验证。

\section{本章小结}

本章系统报告了基于财务合成数据的过程奖励模型训练与评估的完整实验研究。首先介绍了实验环境配置（RTX PRO 6000单卡、Qwen3-8B基座、LoRA微调）与数据集概况（PRM 1,995+400+200，DPO 2800+800+400对），明确了PRM分类性能与置信度校准双维度评估指标以及DPO偏好对齐评估指标。

实验部分通过七个子实验系统验证了PRM训练与DPO偏好优化方案的有效性。实验一至实验四围绕PRM展开：实验一验证了5 epoch训练的充分性（loss下降99.1\%，收敛饱和）；实验二确认步骤准确率达95.25\%/94.00\%（val/test），超过90\%目标；实验三证明了模型置信度校准良好（正确/错误步骤置信度差距0.73）；实验四揭示了小样本PRM训练的过拟合风险——减少epoch（7→5）使错误步骤召回率从0.9\%恢复至79.5\%，是本研究最重要的实验发现。

实验五与实验六聚焦DPO偏好对齐：实验五验证了新版PRM使偏好数据margin从0.004提升至0.348（87倍），并进一步将数据规模扩展至4000对、平均margin提升至0.686（95.4\%样本margin$\geq$0.2），构建了大规模高质量的偏好数据集；实验六先在279对数据上定量验证了数据规模是偏好对齐的瓶颈，随后在2800对数据上完成DPO训练，模型首次成功收敛——loss从随机基线0.693降至0.275、Rewards Margin稳定为正（$\sim$5.0），留出的验证集与测试集偏好准确率分别达99.25\%与99.50\%，远超70\%目标，验证了PRM驱动的DPO偏好对齐路线的有效性。实验七与消融实验进一步系统完成了偏好构造方式（仅步骤平均最优）、$\beta$参数、epoch、打分温度、数据规模等关键组件的消融验证。

%------------------------------------------------------

\newpage
\chapter{结论与展望}
\section{研究总结}

本研究面向银行金融场景下的智能化财务分析需求，针对现有大语言模型在财务推理中存在的数值推理偏差、逻辑连贯性不足等关键瓶颈，系统探索了基于过程监督与偏好对齐的金融推理优化方法。研究遵循"数据合成—过程监督—偏好对齐"的三阶段递进架构，完成了从数据构建到模型训练再到实验验证的完整闭环。

在数据层面，基于Wind/Choice上市公司财务数据库，通过Proposer→Solver→Validator三Agent协作管道，自动合成了1,995条带有步骤级与轨迹级二值标签的金融推理训练样本。每条样本嵌入约70项财务指标作为领域知识，覆盖盈利能力、偿债能力、营运能力、结构性风险和现金流质量五个评估维度。

在模型层面，基于Qwen3-8B基座模型构建了Token分类式过程奖励模型（PRM）。通过在推理轨迹中嵌入步骤分隔符（\texttt{<extra\_0>}）与轨迹结束符（\texttt{<extra\_1>}），配合轻量化2层MLP奖励头和LoRA高效微调（rank=16, $\alpha$=32），实现了对推理步骤质量的细粒度过程监督。实验结果表明，PRM在测试集上达到\textbf{94.00\%的步骤准确率}，正确/错误步骤置信度差距达0.73，具有良好的校准能力。

在偏好对齐层面，基于训练完成的PRM对候选推理轨迹进行评分排序，构建了4000对DPO偏好数据集（平均margin 0.686，95.4\%样本margin$\geq$0.2），并完成了DPO策略模型训练。实验表明，在将训练数据规模从279对扩充至2800对后，DPO训练首次成功收敛——训练loss从随机基线0.693降至0.275，Rewards Margin稳定为正（$\sim$5.0），模型在留出的验证集与测试集上分别取得99.25\%与99.50\%的偏好准确率，远超70\%的预设目标，且无过拟合迹象。这一结果定量验证了"数据规模是偏好对齐的关键瓶颈"的判断，证明PRM驱动的DPO离线偏好对齐路线在金融推理场景下切实可行。

在端到端效果层面，进一步采用大参数模型对DPO策略模型进行三维度评估（结论正确性、步骤正确性、数据引用准确性）。结果表明：策略模型的最终结论正确率较基座提升11.75个百分点（46.00\%至57.75\%），步骤正确率与基座基本持平（61.51\%对62.98\%），数据引用准确性略有下降（30.75\%对35.00\%）。这一结果说明偏好对齐的收益主要集中在最终结论维度，而步骤级过程质量与数据使用准确性仍有提升空间。

在训练策略层面，本研究发现并修复了小样本PRM训练中的严重过拟合问题：旧版配置（7 epoch）下错误步骤召回率仅为0.9\%，模型退化为几乎将所有步骤判为正确。通过将epoch缩减至5，错误步骤召回率恢复至79.5\%，步骤准确率从81.3\%提升至94.0\%。这一发现为小样本PRM训练中的超参数设置提供了重要参考。

\section{创新点与局限性}

本研究的主要创新点可归纳为以下三个方面：

\textbf{（1）大模型合成步骤级二元标注数据}：设计了Proposer→Solver→Validator三Agent协作的数据合成管道，在LangGraph状态图上自动生成包含正负样本的金融推理轨迹，并通过多维度LLM评判（定性评分、过程正确性、事实准确性）为每个推理步骤和轨迹整体标注0/1二值标签。Validator本质上构成了一种LLM双重验证机制：Proposer生成的参考答案与Solver生成的预测答案由同一个235B大模型进行独立比对和验证，确保正负样本标签的有效性和迷惑性。该方法在不依赖人工标注的条件下构建了1,995条高质量的步骤级过程监督数据。

\textbf{（2）Token分类式PRM架构}：在基座模型上引入轻量化奖励头（2层MLP），通过特殊符号标记机制（\texttt{<extra\_0>}和\texttt{<extra\_1>}）将步骤边界显式嵌入输入序列，实现了步骤级与轨迹级双层奖励的端到端联合学习。该架构在保持基座模型语义能力（LoRA仅训练0.75\%参数）的前提下，实现了95\%的步骤准确率和良好校准的置信度输出。

\textbf{（3）PRM驱动的DPO偏好对齐训练}：以过程奖励模型的细粒度评分替代人工偏好标注，自动构建DPO偏好数据集，将传统的GRPO在线奖励搜索范式替换为更简洁稳定的离线偏好对齐流程。实验证明，修复PRM过拟合后，偏好数据的平均区分度从0.004提升至0.686，且在将数据规模扩充至2800对后，DPO策略模型首次收敛，在留出集上取得96\%以上的偏好准确率，验证了该离线偏好对齐路线的有效性，其收益集中体现在最终结论正确性（较基座提升11.75个百分点）的提升上。

本研究同时存在以下局限性：

\textbf{（1）步骤级过程质量未见显著提升}：端到端三维度评估显示，DPO策略模型的步骤正确率与基座基本持平甚至略降（61.51\%对62.98\%），说明步骤级过程监督信号尚未充分传导到策略模型的步骤推理质量，偏好对齐的收益主要集中在最终结论维度。

\textbf{（2）数据引用准确性有待提升}：各模型的数据引用满分率仅约30\%（约60\%样本为部分准确），DPO训练未改善反而略有下降（30.75\%对35.00\%），财务数据使用准确性是当前模型的共性短板。

\textbf{（3）轨迹级准确率有待提升}：当前轨迹准确率为72\%--78\%，低于步骤准确率（94\%--95\%），表明模型在聚合步骤信号推断整体轨迹质量方面仍有改进空间。

\textbf{（4）评分维度消融受限}：三维度评分（定性评分、过程正确性、事实准确性）的独立分数未在数据合成阶段落盘，无法回溯三维度的独立贡献；偏好构造方式、温度、epoch、数据规模等消融实验已完成。

\section{未来研究方向}

基于当前研究成果与局限性分析，后续工作将围绕以下方向展开：

\textbf{（1）步骤级过程质量优化}：探索将步骤级过程监督信号更直接地注入策略模型训练（如步骤级对齐目标），以提升策略模型生成步骤的正确率，弥补当前步骤正确率与基座持平的不足。

\textbf{（2）数据引用准确性提升}：针对财务数据使用准确性设计专门的训练目标或后处理校验机制，提升各模型普遍偏弱（满分率约30\%）的数据引用能力。

\textbf{（3）多epoch训练与奖励黑客风险观察}：当前1 epoch已充分收敛，可尝试2--3 epoch训练，观察是否存在reward hacking风险，进一步验证偏好对齐的稳定性。

\textbf{（4）轨迹准确率优化}：探索步骤信号到轨迹判断的更优聚合策略（如加权投票、注意力聚合等），缩小轨迹级与步骤级准确率之间的差距。

\textbf{（5）多任务与多领域扩展}：将当前财务推理方法推广至反欺诈检测、合规审计等相邻金融任务，验证PRM+DPO训练范式的跨任务泛化能力。


%----------------- 参考文献 -----------------%
\newpage
\Bib{def/GBT7714-2015-NoWarning.bst}{refe.bib}

\vspace{5cm}
\end{document}
